\documentclass[11pt]{article}
\usepackage{tcee}
\addbibresource{3A05.bib}
\tcetema[La síntesis neoclásica. El monetarismo]{3.A.5}{La síntesis neoclásica. El monetarismo.}
\tcefecha{03/11/2026}
\begin{document}
\tceetitulo

\section*{Introducción}
\minutos{2,5}

\begin{itemize}
\item \textbf{\emph{Enganche}}:
  \begin{itemize}
  \item «Cada maestrillo tiene su librillo», dice el refrán. Entre febrero de 2020 y febrero de 2021, la cantidad de dinero en sentido amplio (M2) de Estados Unidos creció casi un 27~\% \parencite{FRED2026M2SL}, y en junio de 2022 la inflación llegó al 9,1~\%, la más alta desde noviembre de 1981 \parencite{BLS2026IPC}. Ante los mismos datos, unos economistas miraban la cantidad de dinero y otros, la brecha de producción y el mercado de trabajo: cada uno con su librillo.
  \item El debate tiene 60 años. En diciembre de 1965 la revista \emph{Time} tituló «ahora todos somos keynesianos» con una frase de \autor{Milton Friedman} \parencite{Time1965Keynesians}, y él aclaró en una carta que la frase completa era «en un sentido, ahora todos somos keynesianos; en otro, nadie es ya keynesiano» \parencite{Friedman1966Time}.
  \item Detrás está la pregunta de esta exposición: \emph{¿qué separaba de verdad a la síntesis neoclásica del monetarismo: el modelo o la forma de leerlo?}
  \end{itemize}

\item \textbf{\emph{Relevancia}}:
  \begin{itemize}
  \item \emph{Para la teoría}: el modelo IS-LM y la curva de \autor{Phillips} aumentada con expectativas siguen siendo el esqueleto de la macroeconomía de los manuales y la primera herramienta para pensar la política fiscal y la monetaria \parencite{DeVroeyHoover2004Introduction}.
  \item \emph{Para la política económica}: los bancos centrales independientes, con un objetivo de inflación y reglas conocidas, son en buena parte herencia del monetarismo \vertema{3.A.36}.
  \item \emph{Para la actualidad}: la inflación de 2021-2023 reabrió la discusión sobre cuánto importa el dinero \parencite{BLS2026IPC}.
  \item Además, como toda historia del pensamiento económico, enseña a no quedarse con el último manual y a ver por qué unas ideas llegan en unas épocas y no en otras: son 2 de los 4 motivos de \textcite{Schumpeter1954History}.\footnote{Los otros 2: la historia es fuente de ideas (principios que no fructificaron pueden orientar la investigación de hoy) y muestra cómo se desplazan las fronteras de la ciencia (cada época trae sus problemas y su forma de afrontarlos). La evolución del pensamiento económico resulta de la interacción de varios factores: el contexto económico de cada momento, el pensamiento anterior y los avances de otras disciplinas, como la filosofía, las matemáticas, la física o la biología. Conocerla permite entender las raíces del pensamiento actual, analizar los fenómenos económicos pasados y presentes y valorar los distintos programas de investigación.}
  \end{itemize}

\item \textbf{\emph{Contextualización}}:
  \begin{itemize}
  \item La historia del pensamiento económico avanza por \emph{revoluciones}: un \emph{paradigma} domina hasta que acumula anomalías que no sabe explicar y otro lo sustituye \parencite{Kuhn1962Structure} \vertema{3.A.1}.
  \item Para los \emph{neoclásicos}, con precios y salarios flexibles, los mercados se vacían: el paro es voluntario o fruto de rigideces y el dinero solo determina el nivel de precios (teoría cuantitativa) \vertema{3.A.3}.
  \item La \emph{Gran Depresión} fue la anomalía: un paro masivo que no se corregía solo. \textcite{Keynes1936} responde con la \emph{Teoría general} \vertema{3.A.4}: la demanda efectiva fija la producción, puede haber paro involuntario sin que nada lo corrija y la política fiscal debe sostener la demanda; bajar los salarios, la receta de \autor{Pigou}, no lo arreglaba.
  \item Pero la \emph{Teoría general} es un libro en prosa y difícil. \autor{Keynes} se dirigía en el prefacio «ante todo a mis colegas economistas» y dejaba en sus manos juzgar si tenía razón \parencite{Keynes1936}. Para \textcite{Mankiw2006Macroeconomist}, leerla hoy es una experiencia «a la vez estimulante y frustrante»: quedan demasiados cabos sueltos y el lector se pregunta cuál es, exactamente, el modelo que une todas las piezas.
  \item La respuesta la dan economistas de distintas tendencias (\autor{John Hicks}, \autor{Alvin Hansen}, \autor{Roy Harrod}, \autor{James Meade}, \autor{Franco Modigliani}, \autor{Paul Samuelson}), que traducen a \autor{Keynes} al lenguaje de los neoclásicos: es la \concepto{síntesis neoclásica}. De los años 40 a finales de los 60 fue la ortodoxia, con el modelo IS-LM y, desde 1960, la curva de \autor{Phillips} \parencite{DeVroeyHoover2004Introduction}. En 1955, \textcite[212]{Samuelson1955Economics} escribía que el 90~\% de los economistas estadounidenses habían dejado de ser keynesianos o antikeynesianos para trabajar en esa síntesis.
  \item Desde mediados de los 60 la inflación de Estados Unidos sube, del 1,6~\% en 1965 a más del 5~\% en 1969 \parencite{BLS2026IPC}, y en los 70 llega la \emph{estanflación}: inflación y paro altos a la vez. La síntesis recibe 3 críticas:
    \begin{itemize}
    \item la de los \emph{postkeynesianos} de Cambridge (\autor{Joan Robinson}), que rechazan el instrumental neoclásico (la controversia del capital);
    \item la de los \emph{keynesianos del desequilibrio} (\autor{Clower}, \autor{Leijonhufvud}, \autor{Malinvaud}), para quienes la síntesis olvida que \autor{Keynes} es una teoría del desequilibrio \vertema{3.A.4};
    \item la \emph{monetarista}, encabezada por \autor{Milton Friedman}: la síntesis había olvidado el dinero, confiaba demasiado en la política discrecional e ignoraba las expectativas. Fue la que más influyó en la política económica, y es la de esta exposición.
    \end{itemize}
  \end{itemize}

\item \textbf{\emph{Problemática}}:
  \begin{itemize}
  \item ¿Cómo se convirtió la \emph{Teoría general} en un modelo con el que hacer política económica?
  \item ¿Qué cambia si el dinero importa y los agentes aprenden de sus errores?
  \item ¿Puede la política de demanda bajar el paro de forma permanente? ¿Reglas o discrecionalidad?
  \item ¿Qué deben hacer la política monetaria, la fiscal, las de oferta y el tipo de cambio si el dinero importa?
  \item ¿Qué queda hoy de cada escuela?
  \end{itemize}

\item \textbf{\emph{Estructura}}:
  \begin{itemize}
  \item La exposición tiene 2 bloques, uno por escuela, en el orden en que se plantearon las preguntas: primero el paradigma y después su crítica.
  \item \concepto{I. La síntesis neoclásica} (\ref{sec:sintesis}): quiénes la formaron y con qué método, su modelo (el IS-LM, del que sale la demanda agregada) y la pieza que lo cerraba, la curva de \autor{Phillips}.
  \item \concepto{II. El monetarismo} (\ref{sec:monetarismo}), con el mismo patrón que la síntesis y algo más:
    \begin{itemize}
    \item la escuela de \autor{Friedman} y su método;
    \item lo que cambia en cada pieza del modelo, primero en la demanda (consumo, inversión y dinero) y después en la oferta (la tasa natural de paro);
    \item sus recomendaciones de política: monetaria y fiscal, de oferta y de economía abierta, con su lectura de la Gran Depresión;
    \item su valoración: lo que queda de él.
    \end{itemize}
  \end{itemize}
\end{itemize}

\begin{esquema}[unbreakable]
\textbf{3.A.5. La síntesis neoclásica. El monetarismo}\\[2pt]
Introducción\esqtiempo{2,5}\\
I. La síntesis neoclásica\esqtiempo{9,5}\\
\hspace*{1.5em}1. Caracterización general de la síntesis neoclásica según \autor{Schumpeter}\\
\hspace*{1.5em}2. El modelo IS-LM (\autor{Hicks}, 1937; \autor{Modigliani}, 1944)\\
\hspace*{1.5em}3. La curva de \autor{Phillips}\\
II. El monetarismo\esqtiempo{17}\\
\hspace*{1.5em}1. Caracterización general de la escuela monetarista según \autor{Schumpeter}\\
\hspace*{1.5em}2. Formalización del pensamiento de \autor{Friedman} y comparación con \autor{Keynes} y la síntesis\\
\hspace*{1.5em}3. Recomendaciones de política económica\\
\hspace*{1.5em}4. Valoración y aspectos criticados (la contracrítica)\\
Conclusión\esqtiempo{1}
\end{esquema}

\begin{notaopositor}
Reparto del cante: introducción, \tiempo{2,5}; I, \tiempo{9,5} (I.1, \tiempo{1,5}; I.2, \tiempo{6}; I.3, \tiempo{2}); II, \tiempo{17} (II.1, \tiempo{2}; II.2, \tiempo{5,5}: demanda, \tiempo{3,5}, y oferta, \tiempo{2}; II.3, \tiempo{7}: política monetaria, \tiempo{1,5}, fiscal, \tiempo{1,5}, expectativas, \tiempo{0,5}, oferta, \tiempo{1}, economía abierta, \tiempo{1,5}, y Gran Depresión, \tiempo{1}; II.4, \tiempo{2,5}), y conclusión, \tiempo{1}. Las cajas grises no se cantan; de la evidencia empírica de la demanda (\ref{sec:demanda}) basta una frase. La síntesis no baja de \tiempo{8} y lleva las fórmulas y los gráficos de la IS, la LM y la demanda agregada: es lo que el tribunal espera ver. Además del refrán del enganche, las 4 pegatinas, cada una con su concepto: «una de cal y otra de arena» (la síntesis), «el hombre es el único animal que tropieza dos veces en la misma piedra» (las expectativas adaptativas), «no se puede engañar a todo el mundo todo el tiempo» (la curva aumentada) y «más vale paso que dure que trote que canse» (la regla del $k$~\%). Los gráficos se construyen por capas: el plano $(Y,i)$ del IS-LM, que vuelve en el debate de las pendientes, y el plano $(U,\pi)$, que empieza con el menú de \autor{Samuelson} y \autor{Solow} y termina con la curva aumentada. Las funciones clave son la IS, la LM y la demanda agregada; después, la curva de \autor{Phillips} aumentada. Las derivaciones algebraicas están en el anexo~\ref{anexo:islm}.
\end{notaopositor}


% ==========================================================================
\section{La síntesis neoclásica}\label{sec:sintesis}
\minutos{9,5}

\begin{itemize}
\item Una vez situada la \emph{Teoría general}, pasamos a ver cómo la convirtió la síntesis neoclásica en un modelo.
\item \emph{Conclusión anticipada}: la síntesis reconcilia a \autor{Keynes} con los neoclásicos. Su modelo, el IS-LM, es keynesiano a corto plazo, por la rigidez del salario nominal, y clásico a largo plazo. Con la curva de \autor{Phillips} queda completo: la política de demanda elige, en un menú, cuánta inflación y cuánto paro.
\end{itemize}

% --------------------------------------------------------------------------
\subsection{Caracterización general de la síntesis neoclásica según \autor{Schumpeter}}\label{sec:caracterizacion-sintesis}

\begin{itemize}
\item \autor{Schumpeter} reconoce una escuela de pensamiento por 3 rasgos: un \emph{líder}, unos \emph{seguidores} y un \emph{cuerpo doctrinal} propio, con su objeto y su método. La síntesis neoclásica los tiene:
  \begin{itemize}
  \item \emph{Líderes}: \autor{John Hicks}, que formula el modelo IS-LM \parencite{Hicks1937Keynes}, y \autor{Paul Samuelson}, que lo lleva a los manuales y llama \concepto{síntesis neoclásica} al consenso que estaba surgiendo en la 3.ª edición de su \emph{Economics} \parencite{Samuelson1955Economics}.
  \item \emph{Seguidores}: \autor{Alvin Hansen}, que lleva a \autor{Keynes} a Estados Unidos desde Harvard; \autor{Franco Modigliani}, \autor{Robert Solow} y el propio \autor{Samuelson} en el MIT; \autor{James Tobin} en Yale; \autor{Lawrence Klein} y la Cowles Commission; \autor{Jan Tinbergen} en los Países Bajos y \autor{James Meade} en el Reino Unido. Casi todos recibieron después el Premio Nobel.\footnote{\autor{Tinbergen} (1969, con \autor{Frisch}), \autor{Samuelson} (1970), \autor{Hicks} (1972, con \autor{Arrow}), \autor{Meade} (1977, con \autor{Ohlin}), \autor{Klein} (1980), \autor{Tobin} (1981), \autor{Modigliani} (1985) y \autor{Solow} (1987).}
  \item \emph{Cuerpo doctrinal}:
    \begin{itemize}
    \item \emph{objeto}: la renta y el empleo a corto plazo y la política de estabilización. Una economía de mercado sujeta a perturbaciones de demanda, sobre todo de la inversión, que la política puede y debe estabilizar;
    \item \emph{método}: modelos de equilibrio con pocas ecuaciones, estimados con la nueva \emph{econometría}. \autor{Tinbergen} construyó en 1936 el primer modelo macroeconométrico de un país, los Países Bajos, y después contrastó las teorías del ciclo \parencite{Tinbergen1939Statistical}; el modelo de \textcite{KleinGoldberger1955Econometric} hizo lo mismo para Estados Unidos. Las expectativas son \emph{estáticas} o exógenas: no reaccionan a la política.
    \end{itemize}
  \end{itemize}
\item \emph{Sus 2 objetivos}:
  \begin{itemize}
  \item \emph{formalizar} a \autor{Keynes}, cuya \emph{Teoría general} apenas tiene matemáticas;
  \item \emph{reconciliarlo} con la teoría neoclásica: \autor{Keynes} para el corto plazo, con salarios y precios rígidos, y los neoclásicos para el largo, con precios flexibles. «Síntesis» quiere decir eso: 2 marcos que se tenían por incompatibles se integran en uno \parencite{DeVroeyDuarte2013Search}.
  \end{itemize}
  Pegatina: «una de cal y otra de arena».
\item \emph{El programa}: perfeccionar cada pieza keynesiana con fundamentos y datos, cada una en su tema:
  \begin{itemize}
  \item \emph{consumo}: la función keynesiana, $C_t=\overline{C}+c\,Y^d_t$, se contrasta con los datos y se completa con la renta relativa y el ciclo vital \vertema{3.A.33};
  \item \emph{inversión}: la eficiencia marginal del capital, $I=\overline{I}-b\,i$, se completa con el acelerador y la $q$ de \autor{Tobin} \vertema{3.A.34};
  \item \emph{dinero}: la preferencia por la liquidez se fundamenta con los modelos de inventario y de cartera \vertema{3.A.35}.
  \end{itemize}
\item \emph{La política}: es la época de la \concepto{sintonía fina} (\emph{fine tuning}), que combina la política fiscal y la monetaria para mantener la economía cerca del pleno empleo. La fiscal era la preferida; la monetaria se tenía por poco eficaz, y en el Reino Unido el Informe \autor{Radcliffe} la relegaba a un papel secundario \parencite{Radcliffe1959Report}.\footnote{Hubo excepciones: \autor{Modigliani} dio desde su modelo de 1944 más peso a la política monetaria que la mayoría de los autores de la síntesis.}
\end{itemize}

% --------------------------------------------------------------------------
\subsection{El modelo IS-LM (\autor{Hicks}, 1937; \autor{Modigliani}, 1944)}\label{sec:islm}

\paragraph{a) Idea}

\begin{itemize}
\item \textcite{Hicks1937Keynes}, en «Mr.~Keynes and the ``Classics''; A Suggested Interpretation», reduce la \emph{Teoría general} y la teoría clásica a un mismo sistema de ecuaciones para ver en qué se diferencian. Llamó LL a su segunda curva. \textcite{Modigliani1944Liquidity} separa los 2 submodelos: el clásico, con salarios flexibles, y el keynesiano, en el que el paro nace de la rigidez a la baja del salario nominal. La etiqueta IS-LM se generalizó con los libros con los que \autor{Alvin Hansen} llevó el modelo a Estados Unidos \parencite{Hansen1949Monetary,Hansen1953Guide}.\footnote{\autor{Harrod}, \autor{Meade} y \autor{Hicks} propusieron en 1936, en la reunión de la Econometric Society de Oxford, sistemas de ecuaciones muy parecidos para resumir la \emph{Teoría general} \parencite{Harrod1937Keynes,Meade1937Simplified,Hicks1937Keynes}. Las ecuaciones ya estaban en \autor{Harrod}, a quien \autor{Keynes} había enviado partes del libro para que las comentara, pero fue el gráfico de \autor{Hicks} el que caló. Por eso, como subraya \textcite{Fonseca2026ISLM}, el IS-LM se lee mejor como la síntesis de varios intentos de formalizar a \autor{Keynes} que como una pieza aislada \parencite{Young1987Interpreting}. El propio \autor{Hicks} se distanció de él al final de su vida: le reprochaba ser un modelo de equilibrio para un problema que, en \autor{Keynes}, es de tiempo y de expectativas \parencite{Hicks1980ISLM}.}
\item Une lo que la cruz keynesiana deja separado, el mercado de bienes y el de dinero, a través del tipo de interés. La pregunta que responde: \emph{¿cómo se determinan a la vez la renta y el tipo de interés a corto plazo, y qué eficacia tienen las políticas fiscal y monetaria?}
\end{itemize}

\paragraph{b) Modelo}

\subparagraph{i) Supuestos}

\begin{itemize}
\item Economía cerrada, a corto plazo, con un nivel de precios dado, $P=\overline{P}$: hay recursos ociosos y la producción se ajusta a la demanda.
\item Sin inflación esperada, $\pi^e=0$, de modo que el tipo de interés nominal y el real coinciden, $i=r$.
\item Tres mercados: bienes, dinero y bonos. Por la ley de \autor{Walras}, si 2 están en equilibrio, el tercero también: basta con estudiar los bienes (la IS) y el dinero (la LM).
\item Las piezas de \autor{Keynes}, en forma lineal:
  \begin{itemize}
  \item consumo, $C=\overline{C}+c\,Y^d$, con $Y^d=(1-t)\,Y$, $0<c<1$ la propensión marginal a consumir y $t$ el tipo impositivo;
  \item inversión, $I=\overline{I}-b\,i$, con $b>0$ su sensibilidad al tipo de interés \vertema{3.A.34};
  \item gasto público exógeno, $\overline{G}$;
  \item demanda de saldos reales por transacción y especulación, $k\,Y-h\,i$, con $k,h>0$ \vertema{3.A.35}, y una oferta monetaria $M$ que fija el banco central.
  \end{itemize}
\end{itemize}

\subparagraph{ii) Desarrollo}

\begin{itemize}
\item \emph{Mercado de bienes: la curva IS}. Es el conjunto de pares $(Y,i)$ con los que la producción iguala la demanda, $Y=C+I+\overline{G}$, o, lo que es lo mismo, la inversión iguala el ahorro (\emph{investment = saving}). Con $\overline{A}\equiv\overline{C}+\overline{I}+\overline{G}$ el gasto autónomo, de $Y=\overline{C}+c\,(1-t)\,Y+\overline{I}-b\,i+\overline{G}$ sale:
  \funcionclave[eq:is]{\text{IS:}\quad Y=\frac{\overline{A}-b\,i}{1-c\,(1-t)},\qquad \overline{A}\equiv\overline{C}+\overline{I}+\overline{G}}
  \begin{itemize}
  \item \emph{En palabras}: un tipo de interés más bajo sube la inversión y, con el multiplicador, la renta. La IS tiene pendiente negativa, $\mathrm{d}i/\mathrm{d}Y=-[1-c\,(1-t)]/b$: es más plana cuanto más responde la inversión al tipo ($b$ grande) y cuanto mayor es el multiplicador.
  \item Una política fiscal expansiva (más $\overline{G}$, menos $t$) o más confianza de los inversores (más $\overline{I}$) la desplazan a la derecha.
  \end{itemize}
\item \emph{Mercado de dinero: la curva LM}. Es el conjunto de pares $(Y,i)$ con los que la demanda de saldos reales iguala la oferta (\emph{liquidity preference = money supply}):
  \funcionclave[eq:lm]{\text{LM:}\quad \frac{M}{P}=k\,Y-h\,i\quad\Longleftrightarrow\quad i=\frac{1}{h}\Big(k\,Y-\frac{M}{P}\Big)}
  \begin{itemize}
  \item \emph{En palabras}: si la renta sube, se quiere más dinero para transacciones; con la oferta dada, el tipo de interés tiene que subir para que se quiera menos por especulación. La LM tiene pendiente positiva, $\mathrm{d}i/\mathrm{d}Y=k/h$: es más plana cuanto más responde la demanda de dinero al tipo ($h$ grande).
  \item Una expansión monetaria (más $M/P$) la desplaza a la derecha.
  \end{itemize}
\item \emph{Equilibrio}: la renta y el tipo de interés se determinan a la vez, donde se cortan las 2 curvas (gráfico~\ref{fig:is-lm}). Sustituyendo la LM en la IS (anexo~\ref{anexo:islm}),
  \begin{equation}
  Y^*=\frac{h\,\overline{A}+b\,M/P}{h\,[1-c\,(1-t)]+b\,k},\qquad i^*=\frac{k\,\overline{A}-[1-c\,(1-t)]\,M/P}{h\,[1-c\,(1-t)]+b\,k}. \label{eq:equilibrio-islm}
  \end{equation}
\end{itemize}

\begin{figure}[H]
\centering
\caption{El modelo IS-LM: equilibrio, expansión fiscal y expansión monetaria}\label{fig:is-lm}
\grafico[0.62\textwidth]{is-lm}
\fuente{elaboración propia a partir de \textcite{Hicks1937Keynes} y \textcite[cap.~1]{Heijdra2009Foundations}.}
\end{figure}

\subparagraph{iii) Implicaciones}

\begin{itemize}
\item \emph{Política fiscal}: más $\overline{G}$ desplaza la IS a la derecha. Al tipo inicial, la renta crecería con el multiplicador simple, de $E$ a $F$; pero con más renta se quiere más dinero, el tipo sube y la inversión cae: es el \concepto{efecto expulsión} (\emph{crowding out}). El nuevo equilibrio, $E'$, tiene más renta y un tipo más alto, y el multiplicador es menor que el keynesiano simple:
  \begin{equation}
  \frac{\partial Y^*}{\partial\overline{G}}=\frac{1}{1-c\,(1-t)+b\,k/h}<\frac{1}{1-c\,(1-t)}. \label{eq:multiplicador-fiscal}
  \end{equation}
\item \emph{Política monetaria}: más $M/P$ desplaza la LM a la derecha; baja el tipo, sube la inversión y, con ella, la renta (de $E$ a $E''$). Es el canal del tipo de interés \vertema{3.A.37}:
  \begin{equation}
  \frac{\partial Y^*}{\partial(M/P)}=\frac{b}{h\,[1-c\,(1-t)]+b\,k}. \label{eq:multiplicador-monetario}
  \end{equation}
\item \emph{Casos extremos}: la eficacia de cada política depende de las pendientes, es decir, de $b$ y de $h$ (tabla~\ref{tab:casos-islm}):
  \begin{itemize}
  \item \concepto{trampa de la liquidez} ($h\to\infty$, LM horizontal): con el tipo tan bajo que todos esperan que suba, el dinero de más se atesora y el tipo no baja. La política monetaria no sirve; la fiscal tiene todo el multiplicador simple. Es el caso de \autor{Keynes} en la depresión;
  \item \emph{caso clásico} ($h=0$, LM vertical): la demanda de dinero no depende del tipo y la LM es la teoría cuantitativa, $Y=(1/k)\,M/P$. La política fiscal expulsa tanta inversión como gasto añade; solo la monetaria mueve la renta;
  \item \emph{inversión insensible al tipo} ($b=0$, IS vertical): la política monetaria baja el tipo, pero no mueve la inversión.
  \end{itemize}
\item Entre los extremos, cuánto sirve cada política depende de cuán inclinadas sean las curvas. Es el \emph{debate de las pendientes}: los keynesianos de la síntesis veían una IS inclinada y una LM plana, y los monetaristas defenderán lo contrario (\ref{sec:demanda}).
\end{itemize}

\begin{table}[H]
\centering
\caption{Los casos extremos del modelo IS-LM}\label{tab:casos-islm}
\small
\begin{tabular}{>{\raggedright\arraybackslash}m{3.7cm}>{\centering\arraybackslash}m{2.3cm}>{\centering\arraybackslash}m{3.4cm}>{\centering\arraybackslash}m{3.4cm}}
\toprule
\rowcolor{tceeazulpalido}\textbf{Caso} & \textbf{Curvas} & \textbf{Multiplicador fiscal} & \textbf{Multiplicador monetario}\\
\midrule
General & IS decreciente, LM creciente & $\dfrac{1}{1-c\,(1-t)+b\,k/h}$ & $\dfrac{b}{h\,[1-c\,(1-t)]+b\,k}$\\
\addlinespace[0.3em]
\textcolor{tceeazul}{\textbf{Trampa de la liquidez}} ($h\to\infty$) & LM horizontal & $\dfrac{1}{1-c\,(1-t)}$ & $0$\\
\addlinespace[0.3em]
\textcolor{tceeverde}{\textbf{Clásico}} ($h=0$) & LM vertical & $0$ & $1/k$\\
\addlinespace[0.3em]
\textcolor{tceeazul}{\textbf{Inversión insensible}} ($b=0$) & IS vertical & $\dfrac{1}{1-c\,(1-t)}$ & $0$\\
\addlinespace[0.3em]
\textcolor{tceeverde}{\textbf{Inversión muy sensible}} ($b\to\infty$) & IS horizontal & $0$ & $1/k$\\
\bottomrule
\end{tabular}
\fuente{elaboración propia a partir de las ecuaciones del modelo (anexo~\ref{anexo:islm}). En azul, los casos keynesianos (solo sirve la política fiscal); en verde, los clásicos (solo sirve la monetaria).}
\end{table}

\begin{itemize}
\item \emph{Con salarios y precios flexibles, ¿hay paro de equilibrio?} Es la pregunta que separa a \autor{Keynes} de los clásicos, y la síntesis la responde en 3 pasos:
  \begin{itemize}
  \item con paro, bajan los salarios y los precios; con $M$ dado, suben los saldos reales $M/P$, la LM se desplaza a la derecha, baja el tipo y sube la inversión hasta el pleno empleo: es el \concepto{efecto \autor{Keynes}};
  \item falla en los 2 casos de \autor{Keynes}, la trampa de la liquidez y la inversión insensible al tipo: ahí la deflación no reactiva la demanda;
  \item pero \textcite{Pigou1943Classical} y \textcite{Patinkin1956Money} añaden que los saldos reales son riqueza: con precios más bajos, el consumo sube aunque el tipo no baje ($\overline{C}$ crece con $M/P$) y la IS se desplaza a la derecha. Es el \concepto{efecto saldos reales} o \concepto{efecto \autor{Pigou}-\autor{Patinkin}} \vertema{3.A.21}: con precios flexibles, la economía siempre vuelve al pleno empleo, aunque el ajuste sea lento \vertema{3.A.4}.
  \end{itemize}
\item \emph{La conclusión de la síntesis}: el paro involuntario existe porque el salario nominal es rígido a la baja \parencite{Modigliani1944Liquidity}. La teoría de \autor{Keynes} pasa a ser un caso particular de la neoclásica, el de los salarios rígidos, y no una «teoría general»; pero el corto plazo era keynesiano y había que estabilizarlo.
\end{itemize}

\paragraph{c) Extensiones}

\begin{itemize}
\item \emph{De la IS-LM a la demanda y la oferta agregadas}. Si el nivel de precios cambia, el equilibrio IS-LM da, para cada $P$, la producción que se demanda: es la \concepto{demanda agregada} (DA),
  \funcionclave[eq:da]{\text{DA:}\quad Y=\frac{h\,\overline{A}+b\,M/P}{h\,[1-c\,(1-t)]+b\,k}}
  \begin{itemize}
  \item \emph{En palabras}: con un nivel de precios más bajo hay más saldos reales, baja el tipo de interés y sube la inversión (el efecto \autor{Keynes}). La DA tiene pendiente negativa en el plano $(Y,P)$ y se desplaza con la política fiscal ($\overline{A}$) y con la monetaria ($M$). En la trampa de la liquidez o con $b=0$ es vertical; el efecto \autor{Pigou}-\autor{Patinkin} la vuelve a inclinar.
  \item \emph{La oferta agregada} (OA) separa el corto plazo del largo (gráfico~\ref{fig:da-oa}):
    \begin{itemize}
    \item a corto plazo, con el salario nominal rígido en $W_0$, es creciente: un precio más alto baja el salario real, $W_0/P$, y las empresas contratan y producen más \parencite{Modigliani1944Liquidity}; con los precios fijos, es horizontal en $\overline{P}$;
    \item a largo plazo, con salarios flexibles, es vertical en la producción natural, $Y^n$, la del pleno empleo del mercado de trabajo.
    \end{itemize}
  \item \emph{El ajuste}: una expansión de demanda sube la producción a corto plazo; con $Y>Y^n$, salarios y precios suben, $M/P$ cae, la LM vuelve a la izquierda y la economía regresa a $Y^n$ con precios más altos. Una expansión monetaria acaba solo en precios (el dinero es neutral a largo plazo) y una fiscal, en un tipo de interés más alto y menos inversión (expulsión total a largo plazo).
  \end{itemize}
\end{itemize}

\begin{figure}[H]
\centering
\caption{De la IS-LM a la demanda y la oferta agregadas: corto y largo plazo}\label{fig:da-oa}
\grafico[0.62\textwidth]{da-oa}
\fuente{elaboración propia a partir de \textcite{Modigliani1944Liquidity} y \textcite[cap.~1]{Heijdra2009Foundations}.}
\end{figure}

\begin{itemize}
\item \emph{Economía abierta}: con la balanza de pagos, el modelo IS-LM-BP de \autor{Mundell} y \autor{Fleming} \vertema{3.A.32}.
\item \emph{La financiación del déficit}: con la restricción presupuestaria del sector público, financiar el gasto con dinero o con deuda no da lo mismo \parencite{BlinderSolow1973Analytical} \vertema{3.A.39}.
\end{itemize}

\paragraph{d) Evidencia empírica}

\begin{itemize}
\item El IS-LM, con las piezas perfeccionadas de la síntesis, fue el esqueleto de los modelos macroeconométricos con los que se hacía la política: el de \textcite{KleinGoldberger1955Econometric} fue uno de los primeros de la economía estadounidense, y en los 60 llegaron modelos de cientos de ecuaciones.
\item Durante unos 25 años después de la Segunda Guerra Mundial, el IS-LM dominó la macroeconomía y estuvo en el centro de los modelos de los gobiernos y de los bancos centrales \parencite{DeVroeyHoover2004Introduction}. Con precios estables y paro bajo, la sintonía fina parecía funcionar.
\end{itemize}

\paragraph{e) Valoración}

\begin{itemize}
\item \emph{Virtudes}:
  \begin{itemize}
  \item \emph{sencillez}: muestra con 2 curvas la interdependencia entre el mercado de bienes y el de dinero;
  \item \emph{versatilidad}: «la versatilidad del modelo es responsable de su supervivencia», decía \textcite{Fischer1986Modiglianiesque}; sirve para la política fiscal y la monetaria, con pleno empleo o con paro, y da resultados cuantitativistas o keynesianos puros según sus pendientes. Por eso sigue en los manuales y en el análisis aplicado.
  \end{itemize}
\item \emph{Críticas}, que abren el resto de la exposición:
  \begin{itemize}
  \item \emph{trabaja con precios dados}: no explica la inflación. La pieza que falta es la curva de \autor{Phillips} (\ref{sec:phillips});
  \item \emph{un solo canal y un solo activo}: el dinero solo actúa a través del tipo de interés de los bonos. Los monetaristas lo verán estrecho (\ref{sec:demanda});
  \item \emph{es estático, con expectativas estáticas y sin microfundamentos}: sus parámetros cambian cuando cambia la política. Es la crítica de \textcite{Lucas1976Econometric} \vertema{3.A.6};
  \item \emph{no explica el paro involuntario como un fallo del mercado}: lo atribuye a un salario rígido que no justifica, y para los keynesianos del desequilibrio se aleja de \autor{Keynes} \vertema{3.A.4}.
  \end{itemize}
\end{itemize}

% --------------------------------------------------------------------------
\subsection{La curva de \autor{Phillips}}\label{sec:phillips}

\begin{itemize}
\item \emph{Idea}: el IS-LM toma el nivel de precios como dado y no explica la inflación, que en los 60 ya no se podía ignorar. La curva de \autor{Phillips} es la pieza que le faltaba \vertema{3.A.28}.
\item \emph{Desarrollo}:
  \begin{itemize}
  \item \textcite{Phillips1958} encuentra en el Reino Unido, entre 1861 y 1957, una relación inversa y convexa entre el ritmo de crecimiento de los \emph{salarios nominales} y la tasa de paro.\footnote{Para \autor{George Akerlof}, la curva de \autor{Phillips} es «probablemente la relación macroeconómica más importante» (traducción propia), en su \href{https://www.nobelprize.org/prizes/economic-sciences/2001/akerlof/lecture/}{conferencia Nobel} de 2001 (minuto 20:32).}
  \item \textcite{SamuelsonSolow1960Analytical} la pasan a \emph{precios}: si las empresas fijan sus precios con un margen fijo sobre el coste laboral unitario, la inflación es el crecimiento de los salarios menos el de la productividad, y en el tramo relevante la curva se aproxima por una recta:
    \begin{equation}
    \pi_t\approx\bar\pi-a\,U_t,\qquad a>0, \label{eq:samuelson-solow}
    \end{equation}
    con $U_t$ la tasa de paro.
  \item Para Estados Unidos presentan un \concepto{menú} de política económica (gráfico~\ref{fig:samuelson-solow}): la estabilidad de precios exigía en torno a un 5,5~\% de paro, y bajarlo al 3~\% costaba una inflación del 4,5~\% anual.
  \end{itemize}
\end{itemize}

\begin{figure}[H]
\centering
\caption{El menú de \autor{Samuelson} y \autor{Solow}: la curva de \autor{Phillips} en precios}\label{fig:samuelson-solow}
\grafico[0.5\textwidth]{samuelson-solow}
\fuente{elaboración propia a partir de \textcite[192, fig.~2]{SamuelsonSolow1960Analytical}.}
\end{figure}

\begin{itemize}
\item \emph{Implicaciones}: con el IS-LM para la demanda y la curva de \autor{Phillips} para los precios, el modelo de la síntesis está completo. La demanda fija la producción y el empleo, y la curva, la inflación; la política de demanda elige en el menú: menos paro a cambio de más inflación, de forma permanente. Desde entonces, todos los modelos macroeconómicos han tenido que explicar esa relación: para \textcite{Akerlof2002Behavioral}, la curva de \autor{Phillips} es «probablemente la relación macroeconómica más importante».
\item \emph{Valoración}:
  \begin{itemize}
  \item los propios \autor{Samuelson} y \autor{Solow} avisaron de que el menú podía cambiar a largo plazo, pero la lectura que se impuso en los 60 fue la de una disyuntiva permanente;
  \item no contempla la inflación de oferta, en la que más inflación llega con más paro, no con menos;
  \item tiene poca fundamentación teórica y no modeliza las expectativas.
  \end{itemize}
  La estanflación de los 70 desplazó la curva hacia arriba, y \autor{Friedman} y \autor{Phelps} explicaron por qué (\ref{sec:oferta}).
\end{itemize}

\begin{itemize}
\item \emph{Recapitulación}: la síntesis dio a la política económica un modelo completo y cuantificable, el IS-LM cerrado con la curva de \autor{Phillips}, y una receta, la sintonía fina. Pero suponía expectativas estáticas y una demanda de dinero inestable, y su menú no resistió la inflación de finales de los 60. Una vez estudiada la síntesis, pasamos a analizar la crítica que la puso en cuestión: la monetarista.
\end{itemize}


% ==========================================================================
\section{El monetarismo}\label{sec:monetarismo}
\minutos{17}

\begin{itemize}
\item Una vez estudiado el paradigma de la síntesis, pasamos a su crítica más influyente: la monetarista.
\item \emph{Conclusión anticipada}:
  \begin{itemize}
  \item El monetarismo no tira el modelo de la síntesis, sino que cambia sus piezas: el consumo y la demanda de dinero son más estables de lo que se creía y el dinero importa más.
  \item Sobre todo, no hay un menú permanente entre inflación y paro: la política de demanda no puede bajar el paro de forma duradera.
  \item De ahí su receta: una regla monetaria sencilla y conocida, una política fiscal sin sintonía fina, reformas de oferta para bajar la tasa natural de paro y tipos de cambio flexibles.
  \end{itemize}
\end{itemize}

% --------------------------------------------------------------------------
\subsection{Caracterización general de la escuela monetarista según \autor{Schumpeter}}\label{sec:caracterizacion-monetarismo}

\begin{itemize}
\item La escuela monetarista cumple los 3 rasgos de \autor{Schumpeter}:
  \begin{itemize}
  \item \emph{Líder}: \autor{Milton Friedman}, profesor en la Universidad de Chicago desde 1946. Desde los años 50 reunió en su taller de dinero y banca a un grupo de discípulos cuyas tesis recogió en \textcite{Friedman1956Quantity}. Recibió el Premio Nobel en 1976 «por sus logros en los campos del análisis del consumo, la historia y la teoría monetarias, y por su demostración de la complejidad de la política de estabilización» \parencite{Nobel1976Friedman}.
  \item \emph{Seguidores}: \autor{Anna Schwartz}, \autor{Phillip Cagan}, \autor{Karl Brunner}, \autor{Allan Meltzer}, \autor{David Laidler}, \autor{Michael Parkin} y, en la economía abierta, \autor{Harry Johnson} \parencite[cap.~4]{SnowdonVane2005Modern}. El nombre de \emph{monetarismo} lo puso en circulación \textcite{Brunner1968Role}.
  \item \emph{Cuerpo doctrinal}, con un objeto y un método comunes.
  \end{itemize}
\item \emph{Objeto}: la macroeconomía y, sobre todo, los efectos de la política económica. Llegan a 3 conclusiones:
  \begin{itemize}
  \item \emph{el dinero importa} (\emph{money matters}): las variaciones de la cantidad de dinero son la causa principal de las de la renta nominal; a corto plazo mueven la producción y a largo plazo solo los precios. Para \textcite{Friedman1963Inflation}, «la inflación es siempre y en todas partes un fenómeno monetario», en el sentido de que solo puede producirla un aumento de la cantidad de dinero más rápido que el de la producción \parencite{Friedman1970Counter} \vertema{3.A.41};
  \item \emph{la economía de mercado es estable}: las grandes perturbaciones vienen de una política monetaria errática, no del sector privado;
  \item \emph{la política discrecional desestabiliza}: mejor reglas. \autor{Friedman} defendió además el libre mercado y un Estado pequeño en \emph{Capitalism and Freedom} \parencite{Friedman1962Capitalism}.
  \end{itemize}
\item \emph{Método}: \textcite{Friedman1953Methodology}, «The Methodology of Positive Economics», en 5 rasgos (desarrollo en el anexo~\ref{anexo:metodo}):
  \begin{enumerate}
  \item \emph{la predicción como criterio}: la economía debe dar predicciones contrastables, y una teoría se mantiene mientras los datos no la refuten; es afín al falsacionismo de \autor{Karl Popper} \vertema{3.A.1};
  \item \emph{economía positiva frente a normativa}: la positiva estudia lo que es, con independencia de los juicios de valor; las discrepancias de política suelen venir de predicciones distintas, no de valores distintos;
  \item \emph{irrelevancia del realismo de los supuestos}: una teoría se juzga por sus predicciones; las empresas se comportan \emph{como si} maximizaran el beneficio, aunque no calculen costes marginales;
  \item \emph{sencillez}: explicar mucho con poco, con el razonamiento económico por delante del aparato técnico;
  \item \emph{enfoque marshalliano}: cada problema por separado, con equilibrio parcial y datos, frente al equilibrio general walrasiano \parencite{Friedman1949Marshallian}.
  \end{enumerate}
\item \emph{Expectativas adaptativas}: la síntesis las trataba como estáticas; los monetaristas las hacen \concepto{adaptativas}. Cada periodo se corrige la previsión en una fracción $\gamma$ del último error. Las usaron \textcite{Cagan1956Monetary} para la inflación esperada en las hiperinflaciones, \textcite{Friedman1957} para la renta permanente \vertema{3.A.33} y \textcite{Nerlove1958Adaptive} para los precios agrícolas de la telaraña, donde amortiguan el ciclo \vertema{3.A.16}. Para la inflación esperada,
  \begin{equation}
  \pi^e_t=\pi^e_{t-1}+\gamma\,\big(\pi_{t-1}-\pi^e_{t-1}\big)=\gamma\sum_{j=0}^{\infty}(1-\gamma)^j\,\pi_{t-1-j},\qquad \gamma\in(0,1]. \label{eq:expectativas-adaptativas}
  \end{equation}
  \begin{itemize}
  \item \emph{En palabras}: la inflación esperada es una media de las pasadas, con pesos que decrecen geométricamente. Con $\gamma=1$ se espera la del periodo anterior, $\pi^e_t=\pi_{t-1}$.\footnote{En la renta permanente, \textcite{Friedman1957} usa la misma regla con la observación del propio periodo, $Y^P_t=Y^P_{t-1}+\gamma\,(Y_t-Y^P_{t-1})$ \vertema{3.A.33}; aquí, como en la telaraña \vertema{3.A.16} y en la curva aceleracionista \vertema{3.A.28}, la previsión se forma con el error del periodo anterior. Es la misma regla con otro calendario. La $\gamma$ sin subíndice es este peso; con subíndice, $\gamma_X$ es la tasa de crecimiento de $X$.}
  \item \emph{Ventaja}: los agentes aprenden de sus errores, lo que las estáticas no permitían.
  \item \emph{Límite}: si la inflación sube año tras año, se quedan siempre cortos y cometen \emph{errores sistemáticos}. Pegatina: «el hombre es el único animal que tropieza dos veces en la misma piedra». \textcite{Muth1961Rational} y la nueva macroeconomía clásica lo resolverán con las expectativas racionales \vertema{3.A.6}.
  \end{itemize}
\end{itemize}

% --------------------------------------------------------------------------
\subsection{Formalización del pensamiento económico de \autor{Friedman}\cortetitulo{} y comparación con \autor{Keynes} y la síntesis neoclásica}\label{sec:formalizacion}

\begin{itemize}
\item La crítica monetarista tiene 2 etapas:
  \begin{itemize}
  \item hasta 1968, se centra en la \emph{demanda}: consumo, inversión y, sobre todo, dinero;
  \item desde \textcite{Friedman1968Role}, en la \emph{oferta}: la tasa natural de paro.
  \end{itemize}
\item Cada pieza se cuenta con el mismo patrón: primero, la visión de \autor{Keynes} y la síntesis; después, la de \autor{Friedman} y los monetaristas, y al final, sus implicaciones.
\end{itemize}

\subsubsection{Crítica por el lado de la demanda}\label{sec:demanda}

\paragraph{a) El consumo: la renta permanente}

\begin{itemize}
\item \emph{\autor{Keynes} y la síntesis}:
  \begin{itemize}
  \item El consumo depende de la renta disponible corriente, $C_t=\overline{C}+c\,Y^d_t$, con $0<c<1$ (la ley psicológica fundamental) y una propensión media que decrece con la renta.\footnote{\autor{Keynes} no supuso que la propensión marginal a consumir fuera constante: «la propensión marginal a consumir no es constante para todos los niveles de empleo, y es probable que, por regla general, tienda a disminuir a medida que el empleo aumenta» \parencite[cap.~10, apdo.~III; traducción propia]{Keynes1936}. La recta con $c$ constante es una simplificación posterior de la síntesis.}
  \item De ahí salen el multiplicador y la eficacia de la política fiscal.
  \item Para \autor{Hansen}, de ahí sale también el riesgo de \emph{estancamiento secular}: si al crecer la renta se ahorra una parte cada vez mayor, la demanda se queda corta \parencite{Hansen1939Economic} \vertema{3.A.40}.
  \end{itemize}
\item \emph{\autor{Friedman}}:
  \begin{itemize}
  \item \textcite{Friedman1957} sustituye la ley psicológica por la \concepto{hipótesis de la renta permanente}. El consumo permanente es una proporción de la renta permanente que no depende de su nivel,
  \begin{equation}
  C^P_t\propto Y^P_t,\qquad Y^P_t=r\,\overline{W}_t, \label{eq:renta-permanente}
  \end{equation}
  con $Y^P_t$ el rendimiento de la riqueza total, humana y no humana, que los hogares estiman con expectativas adaptativas.\footnote{Es la $C^P_t=k\,Y^P_t$ de \textcite{Friedman1957}, con $k$ función del tipo de interés, de la composición de la riqueza y de los gustos \vertema{3.A.33}. Aquí no se escribe esa $k$ para no confundirla con la $k$ de la LM, que mide cuánto responde a la renta la demanda de dinero.}
  \end{itemize}
\item \emph{Implicaciones}:
  \begin{itemize}
  \item Lo transitorio apenas mueve el consumo: una rebaja temporal de impuestos tiene poco efecto y el multiplicador es menor que el keynesiano.
  \item El consumo es más estable que la renta y la propensión media no cae a largo plazo: no hay estancamiento secular y la economía necesita menos estabilización.
  \item La riqueza entra en el consumo, y con ella el dinero y el tipo de interés: es un canal más para la política monetaria.
  \end{itemize}
\end{itemize}

\paragraph{b) La inversión}

\begin{itemize}
\item \emph{\autor{Keynes} y la síntesis}:
  \begin{itemize}
  \item La inversión es el componente inestable de la demanda (los \emph{animal spirits}) y responde poco al tipo de interés \vertema{3.A.34}.
  \item En el IS-LM, $I=\overline{I}-b\,i$ con $b$ pequeño: la IS es inclinada y la política monetaria, débil.
  \end{itemize}
\item \emph{Los monetaristas}:
  \begin{itemize}
  \item El gasto sensible al tipo de interés es mucho más amplio que la inversión empresarial en equipo: incluye la vivienda, los bienes de consumo duradero y las existencias.
  \item Ese gasto responde a toda la gama de rendimientos, no solo al de los bonos.
  \end{itemize}
\item \emph{Implicaciones}: $b$ es grande y la \emph{IS es más plana}. La política monetaria mueve mucho la demanda, y la fiscal expulsa mucho gasto privado.
\end{itemize}

\paragraph{c) La demanda de dinero: la teoría cuantitativa reformulada}

\begin{itemize}
\item \emph{\autor{Keynes} y la síntesis}:
  \begin{itemize}
  \item Se demanda dinero por transacción, precaución y especulación, con solo 2 activos: el dinero y un bono.
  \item La demanda especulativa depende de las expectativas sobre el tipo de interés, que cambian a menudo.
  \item Por eso la demanda de dinero es inestable y muy sensible al tipo ($h$ grande): la LM es plana y la velocidad del dinero, poco fiable \vertema{3.A.35}.
  \end{itemize}
\item \emph{\autor{Friedman}}:
  \begin{itemize}
  \item \textcite{Friedman1956Quantity} reformula la teoría cuantitativa como una teoría de la \emph{demanda} de dinero. El dinero es un activo más de la cartera, y su demanda depende de la riqueza y del rendimiento de las alternativas \parencite[cap.~19]{Mishkin2013Economics}:
  \begin{equation}
  \frac{M^D}{P}=f\Big(\underset{+}{Y^P},\ \underset{-}{r_b-r_m},\ \underset{-}{r_e-r_m},\ \underset{-}{\pi^e-r_m},\ \underset{+}{\overline{W}_H/\overline{W}}\Big), \label{eq:demanda-dinero-friedman}
  \end{equation}
  con $Y^P$ la renta permanente (la riqueza), $r_b$, $r_e$ y $r_m$ los rendimientos esperados de los bonos, las acciones y el propio dinero, $\pi^e$ el de los bienes físicos (su revalorización) y $\overline{W}_H/\overline{W}$ el peso de la riqueza humana, que no se puede vender ni pignorar con facilidad, en la riqueza total.
  \end{itemize}
\item \emph{Implicaciones}:
  \begin{itemize}
  \item \emph{La demanda de dinero es estable}: depende sobre todo de la renta permanente, que cambia despacio, y los diferenciales de rendimiento se mueven poco, porque el del dinero sube con los demás.
  \item \emph{Es poco sensible al tipo de interés} ($h$ pequeño): la \emph{LM es inclinada}, y la trampa de la liquidez, una curiosidad teórica.
  \item \emph{La velocidad es predecible}: si $V$ es estable en $M\,V=P\,Y$, la cantidad de dinero determina la renta nominal.
  \end{itemize}
\end{itemize}

\paragraph{d) El mecanismo de transmisión}

\begin{itemize}
\item \emph{La síntesis}: hay un solo canal, el del tipo de interés, $M\uparrow\ \Rightarrow\ i\downarrow\ \Rightarrow\ I\uparrow\ \Rightarrow\ Y\uparrow$, con 2 eslabones débiles ($h$ grande y $b$ pequeño).
\item \emph{Los monetaristas}: la transmisión es un \emph{reajuste de carteras}.
  \begin{itemize}
  \item Con más dinero del que desean, los agentes compran bonos, acciones, capital ya instalado y bienes duraderos.
  \item Suben los precios de esos activos y bajan sus rendimientos; compensa producir capital nuevo y el gasto sube por muchos caminos a la vez.
  \item \textcite{BrunnerMeltzer1972Money} lo formalizan con 3 activos (dinero, deuda pública y capital real) y un mercado de crédito: el tipo de interés de los bonos es solo uno de los precios relativos que se mueven \parencite{Meltzer1995Monetary} \vertema{3.A.37}.
  \end{itemize}
\end{itemize}

\paragraph{e) Evidencia empírica}

\begin{itemize}
\item \textcite{FriedmanMeiselman1963Relative} compararon en Estados Unidos (1897-1958) el poder explicativo de la cantidad de dinero y del gasto autónomo del multiplicador: el consumo se explicaba mejor con el dinero.
\item La réplica de \textcite{AndoModigliani1965Relative}, con otra definición del gasto autónomo, daba el resultado contrario: la evidencia dependía de cómo se midieran las variables, y el debate siguió abierto.
\end{itemize}

\paragraph{f) El debate de las pendientes}

\begin{itemize}
\item En el IS-LM, el debate de los años 60 se resumió en las pendientes de las 2 curvas (gráfico~\ref{fig:pendientes}):
  \begin{itemize}
  \item \emph{keynesianos}: IS inclinada ($b$ pequeño) y LM plana ($h$ grande): la política fiscal es eficaz, con poca expulsión, y la monetaria, débil;
  \item \emph{monetaristas}: IS plana ($b$ grande) y LM inclinada ($h$ pequeño): la política monetaria es potente y la fiscal se expulsa a sí misma.
  \end{itemize}
\item \emph{Matiz}: \autor{Friedman} nunca aceptó que su tesis dependiera de una LM vertical, y solo usó el IS-LM, en 1970, para hablar el lenguaje de sus críticos. Lo que le separaba de ellos no era una pendiente, sino la estabilidad de la demanda de dinero y la amplitud de la transmisión; él mismo lo explicaba por el método: él era marshalliano y sus críticos, walrasianos \parencite{BordoSchwartz2003ISLM}.
\end{itemize}

\begin{figure}[H]
\centering
\caption{El debate de las pendientes: la misma expansión fiscal y monetaria según keynesianos y monetaristas}\label{fig:pendientes}
\grafico[0.9\textwidth]{pendientes}
\fuente{elaboración propia a partir de \textcite{BordoSchwartz2003ISLM} y \textcite[cap.~1]{Heijdra2009Foundations}.}
\end{figure}

\subsubsection{Crítica por el lado de la oferta}\label{sec:oferta}

\begin{itemize}
\item \emph{Idea}: la síntesis tomaba la curva de \autor{Phillips} como un menú estable. En su discurso presidencial ante la American Economic Association, \textcite{Friedman1968Role}, y a la vez \textcite{Phelps1967Phillips,Phelps1968Money}, explican que a los trabajadores y a las empresas les importa el salario \emph{real}, y que al negociar el nominal prevén la inflación \vertema{3.A.28}.
\item \emph{Modelo}: la curva de \autor{Phillips} \concepto{aumentada con expectativas},
  \funcionclave[eq:phillips-aumentada]{\pi_t=\pi^e_t-a\,\big(U_t-U^n\big)+\varepsilon_t,\qquad a>0}
  con $\pi^e_t\equiv\mathbb{E}_{t-1}[\pi_t]$ la inflación esperada y $\varepsilon_t$ las perturbaciones de oferta.
  \begin{itemize}
  \item La \concepto{tasa natural de paro}, $U^n$, es la que hay cuando la inflación es la esperada. Para \autor{Friedman}, es «el nivel que resultaría del sistema walrasiano de ecuaciones de equilibrio general, siempre que incorpore las características estructurales reales de los mercados de trabajo y de bienes» \parencite[8]{Friedman1968Role}. La fijan factores reales, no la demanda \vertema{3.A.26}.
  \item \emph{Hay una curva de corto plazo para cada inflación esperada} (gráfico~\ref{fig:phillips-aumentada}): una expansión de demanda sube los precios; con $\pi^e$ dada, el salario real cae, se contrata más y el paro baja (de $A$ a $B$). Cuando los trabajadores revisan sus expectativas, el paro vuelve a $U^n$ con más inflación ($C$). A largo plazo, con $\pi=\pi^e$, la curva es \emph{vertical} en $U^n$. Pegatina: «no se puede engañar a todo el mundo todo el tiempo».
  \item Con $\pi^e_t=\pi_{t-1}$ ($\gamma=1$ en \eqref{eq:expectativas-adaptativas}) y sin perturbaciones de oferta, es la \concepto{curva aceleracionista}:
    \begin{equation}
    \Delta\pi_t\equiv\pi_t-\pi_{t-1}=-a\,\big(U_t-U^n\big). \label{eq:aceleracionista}
    \end{equation}
    Mantener el paro por debajo de $U^n$ exige una inflación cada vez mayor (de $C$ a $D$). De ahí el nombre de NAIRU, la tasa de paro que no acelera la inflación; su versión con fijación de salarios y precios es posterior \vertema{3.A.28}.
  \end{itemize}
\end{itemize}

\begin{figure}[H]
\centering
\caption{La curva de \autor{Phillips} aumentada con expectativas: corto y largo plazo}\label{fig:phillips-aumentada}
\grafico[0.6\textwidth]{phillips-aumentada}
\fuente{elaboración propia a partir de \textcite{Friedman1968Role} y \textcite{Phelps1967Phillips}.}
\end{figure}

\begin{itemize}
\item \emph{Implicaciones}: «siempre hay una disyuntiva temporal entre inflación y paro; no hay una disyuntiva permanente» \parencite[11]{Friedman1968Role}. La política de demanda no puede fijar el paro; para bajar $U^n$ hacen falta políticas de oferta.
\item \emph{Evidencia}: la predicción llegó antes que los datos. En Estados Unidos, entre 1969 y 1971 el paro subió del 3,5~\% al 6~\% y la inflación apenas bajó, del 5,4~\% al 4,2~\%; en 1975, tras el choque del petróleo, convivían un paro del 8,5~\% y una inflación del 9,1~\%: la \emph{estanflación} \parencite{BLS2026UNRATE,BLS2026CPI}. Las curvas de corto plazo se habían desplazado hacia arriba con la inflación esperada.
\item \emph{Valoración}: separa el corto plazo del largo y pone las expectativas en el centro de la macroeconomía. Pero, con expectativas adaptativas, los agentes se equivocan siempre en la misma dirección mientras la inflación sube: es la puerta de la nueva macroeconomía clásica \parencite{Lucas1972Expectations} \vertema{3.A.6}.
\end{itemize}

\begin{itemize}
\item \emph{Recapitulación de la formalización}: los monetaristas dejan en pie el IS-LM, pero con otras pendientes y un dinero más potente, y lo cierran con otra curva de \autor{Phillips}, vertical a largo plazo. Visto qué cambia en el modelo, pasamos a ver qué cambia en la política económica.
\end{itemize}

% --------------------------------------------------------------------------
\subsection{Recomendaciones de política económica}\label{sec:politica}

\begin{itemize}
\item \emph{Idea}: las recomendaciones monetaristas salen de las 3 piezas de la formalización, y cada una cierra una puerta a la sintonía fina de la síntesis:
  \begin{itemize}
  \item \emph{el dinero importa}: la política monetaria es potente, pero, mal llevada, es la principal fuente de perturbaciones;
  \item \emph{la demanda privada es estable}: no hace falta afinar la coyuntura, y quien lo intenta llega tarde por los retardos;
  \item \emph{la tasa natural y las expectativas adaptativas}: la política de demanda solo baja el paro mientras engaña a los agentes; lo que lo baja de forma duradera son las políticas de oferta.
  \end{itemize}
\item Las recomendaciones se ordenan con el mismo patrón que la formalización (primero la síntesis, después los monetaristas): políticas de demanda, separando la monetaria y la fiscal; políticas de oferta, y economía abierta. Cierra el apartado su lectura de la Gran Depresión, la prueba de cargo de todo lo anterior.
\end{itemize}

\subsubsection{Políticas de demanda}\label{sec:politica-demanda}

\paragraph{a) Política monetaria}\label{sec:politica-monetaria}

\begin{itemize}
\item \emph{\autor{Keynes} y la síntesis}:
  \begin{itemize}
  \item La política monetaria es un instrumento secundario, débil por la LM plana y la IS inclinada, y casi inútil en una depresión, como en la trampa de la liquidez.
  \item Si se usa, es para fijar el tipo de interés y las condiciones del crédito, y de forma discrecional, junto a la fiscal, para afinar la coyuntura. En el Reino Unido, el Informe \autor{Radcliffe} la dejaba en un papel secundario \parencite{Radcliffe1959Report}.
  \end{itemize}
\item \emph{\autor{Friedman} y los monetaristas}:
  \begin{itemize}
  \item \emph{Es potente, pero actúa con retardos largos y variables} \parencite{Friedman1961Lag}:
    \begin{itemize}
    \item los \emph{retardos internos} van de que surge el problema a que se reconoce y se decide actuar;
    \item los \emph{retardos externos} van de la decisión a que hace efecto en la renta, y son largos y cambian de un ciclo a otro;
    \item cuando la medida hace efecto, la situación ha cambiado, y la sintonía fina acaba amplificando el ciclo que pretende suavizar.
    \end{itemize}
  \item \emph{No se sabe de antemano cuánto del efecto irá a la producción y cuánto a los precios}: \autor{Friedman} reconocía que a su marco le faltaba esa ecuación, la que reparte la variación de la renta nominal entre precios y cantidades \parencite{Friedman1970Theoretical}. Otra razón para no afinar.
  \item \emph{Lo que la política monetaria puede y no puede hacer} \parencite{Friedman1968Role}:
    \begin{itemize}
    \item \emph{no puede} fijar el tipo de interés ni la tasa de paro más que por poco tiempo: si lo intenta, acaba en una inflación (o una deflación) creciente;
    \item \emph{sí puede} impedir que el dinero sea una fuente de perturbaciones, dar un marco estable a las decisiones privadas (un nivel de precios predecible) y compensar perturbaciones grandes y claras.
    \end{itemize}
  \item \emph{El objetivo intermedio es la cantidad de dinero, no el tipo de interés}:
    \begin{itemize}
    \item una expansión monetaria baja el tipo de interés al principio, pero después la mayor renta y, sobre todo, la mayor inflación esperada lo suben (el efecto \autor{Fisher}) \parencite{Friedman1968Role};
    \item por eso los tipos altos pueden ser señal de una política que ha sido laxa, y los bajos, de una que ha sido restrictiva: el tipo de interés es un mal indicador de la política;
    \item el banco central controla la cantidad de dinero a través de la base monetaria, porque el multiplicador es estable; \autor{Brunner} y \autor{Meltzer} defendían controlar la base \parencite[cap.~4]{SnowdonVane2005Modern} \vertema{3.A.36}.
    \end{itemize}
  \item \emph{La regla del $k$~\%}: \textcite{Friedman1960Program} propone que la cantidad de dinero crezca a un ritmo fijo, año tras año, sin responder al ciclo. Con la ecuación cuantitativa en tasas,
  \begin{equation}
  \gamma_M+\gamma_V=\pi+\gamma_Y, \label{eq:cuantitativa}
  \end{equation}
  si la velocidad es estable, un crecimiento del dinero igual al de la producción a largo plazo da precios estables. La regla es
  \begin{equation}
  \gamma_M=\overline{\gamma}_M, \label{eq:regla-k}
  \end{equation}
  con $\overline{\gamma}_M$ constante, que \autor{Friedman} situaba entre el 3 y el 5~\% anual: el crecimiento de la producción más un margen por la caída tendencial de la velocidad.
    \begin{itemize}
    \item \emph{Ventajas}: no depende de previsiones, no da margen a la presión política y ancla las expectativas \vertema{3.A.36}.
    \item \emph{Es un estabilizador automático}: en una recesión, el dinero sigue creciendo al mismo ritmo y su crecimiento supera al de la renta; en un auge, se queda por debajo.
    \item Pegatina: «más vale paso que dure que trote que canse».\footnote{No hay que confundirla con la \emph{regla de \autor{Friedman}} de la cantidad óptima de dinero \parencite{Friedman1969Optimum}: como producir dinero no cuesta casi nada, su coste de oportunidad, el tipo de interés nominal, debería ser cero, lo que exige una deflación igual al tipo de interés real \vertema{3.A.41}.}
    \end{itemize}
  \end{itemize}
\end{itemize}

\begin{anotaciones}
\textbf{La regla en detalle (no se canta)}
\begin{itemize}
\item \emph{Reglas frente a autoridades}: \autor{Friedman} quería que la regla la fijara la ley, para atar las manos del banco central y quitarle la tentación de actuar en cada coyuntura \parencite{Friedman1960Program,Friedman1962Capitalism}. En la misma obra de 1960 proponía un coeficiente de reservas del 100~\% sobre los depósitos, para que la cantidad de dinero dependiera solo de la base que crea el banco central \parencite{Friedman1960Program}.
\end{itemize}
\end{anotaciones}

\begin{itemize}
\item \emph{La regla en la práctica: el experimento monetarista}:
  \begin{itemize}
  \item En Estados Unidos, el 6 de octubre de 1979, la Reserva Federal de \autor{Paul Volcker} pasó a controlar las reservas bancarias para frenar el crecimiento del dinero y dejó que el tipo de interés se moviera libremente \parencite{Medley2013Volcker}.
  \item La inflación cayó a menos de la tercera parte en 3 años, a costa del paro más alto de la posguerra hasta entonces (cifras en la caja).
  \item En el Reino Unido, la estrategia financiera a medio plazo (\emph{Medium-Term Financial Strategy}) del presupuesto del 26 de marzo de 1980 fijó una senda decreciente de objetivos para la oferta monetaria amplia (£M3), empezando por un crecimiento del 7-11~\% anual \parencite{Howe1980Budget}.
  \item Pero la velocidad dejó de ser estable (gráfico~\ref{fig:velocidad}): con la desregulación financiera y los nuevos depósitos remunerados, la relación entre dinero y renta se rompió, y los bancos centrales abandonaron uno tras otro los objetivos monetarios.
  \end{itemize}
\end{itemize}

\begin{anotaciones}
\textbf{El experimento monetarista en cifras (no se canta)}
\begin{itemize}
\item \emph{Estados Unidos}: el tipo de los fondos federales llegó al 19,10~\% de media en junio de 1981 \parencite{FRED2026FondosFederales}; la inflación bajó del 14,8~\% de marzo de 1980 al 3,8~\% de diciembre de 1982, con un paro del 10,8~\% en noviembre de 1982 \parencite{BLS2026IPC,BLS2026UNRATE}. En octubre de 1982 la Reserva Federal dejó de fijar un objetivo para M1 \parencite{Thornton2006When}.
\item \emph{La velocidad}: la de M1, que crecía de forma regular, rompió su tendencia a partir de 1981; la de M2, estable hasta 1990, subió hasta 2,19 en 1997 y cayó a 1,13 en 2020, y hoy es de 1,42 \parencite{FRED2026Velocidad}.
\item \emph{El abandono de los objetivos monetarios}: el Banco de Canadá abandonó sus objetivos de M1 en noviembre de 1982 \parencite{Bouey1982Statement}; en julio de 1993, la Reserva Federal dijo al Congreso que M2 había dejado de ser un indicador fiable \parencite{Greenspan1993Testimony}, y en 2000 dejó de fijar bandas para los agregados monetarios \parencite{Fed2000MPR}.
\end{itemize}
\end{anotaciones}

\begin{figure}[H]
\centering
\caption{Velocidad del dinero en Estados Unidos, 1959-2026 (PIB nominal entre M2 y entre M1)}\label{fig:velocidad}
\grafico[0.85\textwidth]{velocidad-dinero}
\fuente{elaboración propia a partir de \textcite{FRED2026Velocidad}; último dato, tercer trimestre de 2026. La velocidad de M1 se corta en 2020 porque la Reserva Federal cambió la definición de M1 en mayo de ese año.}
\end{figure}

\paragraph{b) Política fiscal}\label{sec:politica-fiscal}

\begin{itemize}
\item \emph{\autor{Keynes} y la síntesis}:
  \begin{itemize}
  \item La política fiscal es el instrumento preferido: con la IS inclinada y la LM plana, el multiplicador es grande y la expulsión, pequeña \eqref{eq:multiplicador-fiscal}.
  \item Se usa de forma discrecional, con déficits en las recesiones, y la renta corriente mueve el consumo, así que también sirven las rebajas temporales de impuestos.
  \end{itemize}
\item \emph{\autor{Friedman} y los monetaristas}:
  \begin{itemize}
  \item \emph{La expulsión es grande}: con la IS plana y la LM inclinada, un gasto público financiado con deuda sube el tipo de interés y desplaza mucho gasto privado sensible al tipo. A largo plazo, con la producción en su nivel natural, la expulsión es total (tabla~\ref{tab:casos-islm}, caso clásico): el gasto público cambia el reparto de la producción entre el sector público y el privado, no su nivel \parencite{Friedman1970Counter}.
  \item \emph{Las rebajas transitorias de impuestos apenas mueven el consumo}, por la renta permanente (apartado~\ref{sec:demanda}).
  \item \emph{Lo que cuenta es cómo se financia el déficit}: la política fiscal «pura», financiada con deuda, mueve poco la renta nominal; la que la mueve es la que va acompañada de dinero, y eso es política monetaria. \autor{Friedman} lo debatió en público con \autor{Walter Heller}, que había presidido el Consejo de Asesores Económicos de \autor{Kennedy} \parencite{FriedmanHeller1969Monetary} \vertema{3.A.38}.
  \item \emph{La financiación monetaria del déficit es inflacionista}: si los déficits se cubren con dinero de forma continuada, la inflación sube y, con expectativas que se adaptan cada vez más deprisa, puede acabar en hiperinflación, como mostró \textcite{Cagan1956Monetary} \vertema{3.A.39}.
  \item \emph{Reglas también para el presupuesto}: \textcite{Friedman1948Framework} proponía no tocar el gasto ni los tipos impositivos por la coyuntura, fijarlos por lo que la sociedad quiere a largo plazo, equilibrar el presupuesto en pleno empleo y dejar que los \emph{estabilizadores automáticos} (impuestos y prestaciones que se mueven solos con el ciclo) dieran déficits en las recesiones y superávits en los auges.
  \item \emph{Un Estado pequeño}: \textcite{Friedman1962Capitalism} propone un impuesto negativo sobre la renta en lugar de los programas de asistencia y cheques escolares en lugar de la provisión pública.\footnote{Un ejemplo de los cheques escolares, el sistema educativo sueco, en \href{https://youtu.be/DkGyH-E82uI?t=146}{este vídeo de VisualEconomik} (YouTube).}
  \end{itemize}
\end{itemize}

\begin{anotaciones}
\textbf{La financiación del déficit y la expulsión (no se canta)}
\begin{itemize}
\item \emph{Más canales de expulsión}: si los tenedores de deuda pública la ven como riqueza, la emisión de bonos sube también la demanda de dinero (la LM se desplaza a la izquierda) y la expulsión puede ser mayor que la que dan solo las pendientes; con la restricción presupuestaria del sector público, financiar con dinero o con deuda no da lo mismo \parencite{BlinderSolow1973Analytical}.
\item \emph{En el marco de 1948}, los déficits de las recesiones se financiaban creando dinero y los superávits de los auges lo retiraban: la política monetaria y la fiscal se coordinaban con una regla, sin decisiones discrecionales \parencite{Friedman1948Framework}. En 1960, \autor{Friedman} lo cambió por la regla del $k$~\%, más sencilla \parencite{Friedman1960Program}.
\item \emph{Extensión}: \textcite{SargentWallace1981} mostraron que, si la política fiscal no se ajusta, el banco central monetarista acaba teniendo que financiar el déficit, antes o después, y el control monetario de hoy puede traer más inflación mañana (la dominación fiscal) \vertema{3.A.39}.
\end{itemize}
\end{anotaciones}

\paragraph{c) Por qué no hay atajos: las expectativas y la tasa natural}\label{sec:politica-expectativas}

\begin{itemize}
\item \emph{La política de demanda solo baja el paro mientras engaña}: con expectativas adaptativas, una expansión baja el paro mientras la inflación supera a la esperada. Mantener el paro por debajo de $U^n$ exige una inflación cada vez mayor \eqref{eq:aceleracionista}: es la \emph{tesis aceleracionista}.
\item \emph{Por eso el objetivo de la política de demanda es la estabilidad nominal}: un crecimiento estable del dinero y, con él, de la renta nominal y de los precios, no una tasa de paro.
\item \emph{La desinflación tiene un coste}: con expectativas adaptativas, bajar la inflación exige mantener un tiempo el paro por encima de $U^n$ \vertema{3.A.41}.
  \begin{itemize}
  \item Los monetaristas preferían el \emph{gradualismo}: reducir el crecimiento del dinero poco a poco y anunciarlo, para que las expectativas se ajustaran antes.
  \item Para abaratarla, \textcite{Friedman1974Monetary} proponía \emph{indexar} los contratos (cláusulas de revisión según la inflación): así, los salarios y los tipos de interés bajan solos con la inflación y el paro sube menos.
  \item La nueva macroeconomía clásica, con expectativas racionales, defenderá en cambio una desinflación de golpe y creíble \vertema{3.A.6}.
  \end{itemize}
\end{itemize}

\subsubsection{Políticas de oferta}\label{sec:politica-oferta}

\begin{itemize}
\item \emph{\autor{Keynes} y la síntesis}:
  \begin{itemize}
  \item La oferta casi no se trata: la producción la fija la demanda y la curva de \autor{Phillips} es un menú dado.
  \item Para mejorar el menú, la síntesis recurre a las \emph{políticas de rentas}: pautas o controles de salarios y precios que frenen la inflación sin subir el paro.
  \end{itemize}
\item \emph{\autor{Friedman} y los monetaristas}:
  \begin{itemize}
  \item \emph{Solo las reformas bajan la tasa natural}: la fijan las características estructurales de los mercados, y muchas las crean las leyes y la política. El salario mínimo y el poder de los sindicatos la suben; mejores servicios de empleo y más información sobre las vacantes y la oferta de trabajo la bajan \parencite[9]{Friedman1968Role} \vertema{3.A.28}.
  \item \emph{No a las políticas de rentas}: si la inflación es un fenómeno monetario, los controles de precios y salarios no la curan, solo la reprimen mientras duran y distorsionan los precios relativos \parencite[cap.~4]{SnowdonVane2005Modern}.
  \item \emph{Más mercado}: menos trabas a la contratación y a la movilidad, menos poder de monopolio de empresas y sindicatos y un impuesto negativo sobre la renta que no desanime a trabajar como los subsidios tradicionales \parencite{Friedman1962Capitalism}.
  \item \emph{La estabilidad de precios es también una política de oferta}: en su conferencia del Nobel, \textcite{Friedman1977Nobel} sostiene que una inflación alta y variable puede subir el paro, porque enturbia las señales de los precios relativos y lleva a instituciones y contratos peor adaptados. A medio plazo, la curva de \autor{Phillips} puede tener pendiente positiva.
  \end{itemize}
\end{itemize}

\subsubsection{Economía abierta}\label{sec:economia-abierta}

\begin{itemize}
\item \emph{\autor{Keynes} y la síntesis}:
  \begin{itemize}
  \item Defienden los tipos de cambio fijos pero ajustables de Bretton Woods y tratan la balanza de pagos como un problema de flujos de bienes: el ajuste viene por las elasticidades de las exportaciones e importaciones y por la absorción, con devaluaciones y, si hace falta, controles \vertema{3.B.12}.
  \item En el modelo IS-LM-BP de \autor{Mundell} y \autor{Fleming}, la eficacia de cada política depende del régimen cambiario y de la movilidad del capital \vertema{3.A.32}.
  \end{itemize}
\item \emph{\autor{Friedman}: los tipos de cambio flexibles}. \textcite{Friedman1953Case} los defiende cuando regía Bretton Woods, con 4 argumentos:
  \begin{itemize}
  \item \emph{ajustan los precios relativos entre países sin deflación ni paro}: con precios y salarios rígidos, es más fácil mover un solo precio, el tipo de cambio, que todos los precios internos, igual que es más fácil adelantar la hora que cambiar los horarios de todo el mundo;
  \item \emph{dejan a cada país elegir su política monetaria}, y con ella su regla y su inflación: con tipo fijo, la inflación de fuera se importa;
  \item \emph{la especulación es estabilizadora}: quien especula contra los fundamentos compra caro y vende barato, pierde dinero y desaparece;
  \item \emph{hacen innecesarios los controles} de cambios y de comercio con que se defendían las paridades.
  \end{itemize}
  El fin de Bretton Woods (1971-1973) llevó a la flotación generalizada \vertema{3.B.20} \vertema{3.B.15}.
\item \emph{\autor{Harry Johnson}: el enfoque monetario de la balanza de pagos}. Desde Chicago y la London School of Economics, con \autor{Robert Mundell} y \autor{Jacob Frenkel}, trata la balanza de pagos como un fenómeno monetario \parencite{FrenkelJohnson1976Monetary}:
  \begin{itemize}
  \item \emph{Supuestos}: demanda de dinero estable, $M^D=k\,P\,Y$; producción en su nivel natural; paridad del poder adquisitivo, $P=e\,P^*$; y una oferta monetaria que es la suma del crédito interno, $D$, y de las reservas internacionales, $R$: $M^S=D+R$.
  \item \emph{Resultado}: con tipo de cambio fijo, las reservas se ajustan hasta que la oferta de dinero iguala la demanda,
  \begin{equation}
  \Delta R=\Delta M^D-\Delta D. \label{eq:enfoque-monetario}
  \end{equation}
  Un déficit exterior es un exceso de oferta de dinero: si el banco central crea más crédito interno del que se demanda, el exceso se va al exterior y se pierden reservas.
  \item \emph{Implicaciones}:
    \begin{itemize}
    \item los desequilibrios de la balanza de pagos son transitorios y se corrigen solos con la variación de las reservas, si el banco central no las esteriliza;
    \item una devaluación solo mejora la balanza mientras sube $P$ y con él la demanda de dinero: su efecto es transitorio;
    \item un país que crece más gana reservas, porque se demanda más dinero;
    \item con tipo de cambio flexible, el mismo razonamiento da el \emph{enfoque monetario del tipo de cambio}: la moneda del país que crea más dinero se deprecia \vertema{3.B.14}.
    \end{itemize}
  \end{itemize}
\end{itemize}

\subsubsection{Teoría de los ciclos: la Gran Depresión}\label{sec:ciclos}

\begin{itemize}
\item \emph{\autor{Keynes} y la síntesis}: la Gran Depresión mostraba que la demanda privada puede hundirse y quedarse hundida; como en la trampa de la liquidez, la política monetaria no servía y la salida era fiscal.
\item \emph{\autor{Friedman} y \autor{Schwartz}}: en el capítulo~7 de \emph{A Monetary History of the United States, 1867-1960} \parencite{FriedmanSchwartz1963Monetary} dan la vuelta al argumento con los datos. La Gran Contracción de 1929-1933 no prueba la impotencia de la política monetaria, sino su poder: es «un trágico testimonio de la importancia de las fuerzas monetarias» \parencite{FriedmanSchwartz1965Great}.
  \begin{itemize}
  \item \emph{Los hechos} (gráfico~\ref{fig:gran-contraccion}): del pico de agosto de 1929 al fondo de marzo de 1933, la cantidad de dinero cayó un 35~\%, de 46.280 a 29.970 millones de dólares. Más de una quinta parte de los bancos comerciales suspendió pagos y, con fusiones y liquidaciones, su número pasó de 24.504 en junio de 1929 a 13.949 en junio de 1933 \parencite{FriedmanSchwartz1963Monetary,FRED2026Bancos}. Entre 1929 y 1933, el PIB real cayó un 26~\% y los precios de consumo, un 25~\% \parencite{BEA2026NIPA,BLS2026IPC}.\footnote{\autor{Friedman} y \autor{Schwartz}, con las cuentas nacionales de los años 60, daban una caída de la renta real del 36~\%. Las revisiones posteriores de la contabilidad nacional la han suavizado; aquí se usan los datos actuales del BEA.}
  \item \emph{El mecanismo}: la base monetaria no cayó; por sí sola habría elevado la cantidad de dinero un 17,5~\%. Lo que la hundió fue el multiplicador: con las quiebras, el público quiso más efectivo y menos depósitos (el cociente depósitos/efectivo bajó de 10,8 a 4,4) y los bancos, más reservas (el cociente depósitos/reservas bajó de 13,1 a 9,1) \parencite{FriedmanSchwartz1965Great} \vertema{3.A.36}.
  \item \emph{La culpa}: la Reserva Federal no compró títulos ni prestó la liquidez que habría frenado los pánicos bancarios. Si hubiera actuado a finales de 1930 o en 1931, la contracción habría sido menos grave y, muy probablemente, habría acabado «mucho antes» \parencite{FriedmanSchwartz1965Great}.\footnote{La contracción no vino precedida de una expansión monetaria: el dinero llevaba casi 2 años sin crecer, porque la Reserva Federal había endurecido su política en 1928-1929 para frenar la especulación bursátil \parencite{FriedmanSchwartz1965Great}. La tesis de que una política monetaria laxa en los años 20 infló una burbuja que estalló en 1929 es de la escuela austriaca \parencite{Rothbard1963America}, no de \autor{Friedman} y \autor{Schwartz}.}
  \end{itemize}
\end{itemize}

\begin{figure}[H]
\centering
\caption{La Gran Contracción en Estados Unidos, 1929-1933: dinero, precios, producción y los cocientes del multiplicador}\label{fig:gran-contraccion}
\grafico[0.95\textwidth]{gran-contraccion}
\fuente{elaboración propia a partir de \textcite{NBER2026Macrohistory} (series de \autor{Friedman} y \autor{Schwartz}), \textcite{BLS2026IPC}, \textcite{FRED2026ProduccionIndustrial} y \textcite{FRED2026Suspensiones}. La serie de bancos suspendidos acaba en febrero de 1933. Variaciones del gráfico de pico a fondo con datos mensuales; en el texto, medias anuales de 1929-1933 y PIB real.}
\end{figure}

\begin{itemize}
\item \emph{La lección aprendida}: \autor{Ben Bernanke} estudió la crisis financiera de los años 30 siguiendo a \autor{Friedman} y \autor{Schwartz}, y añadió el canal del crédito: la quiebra de los bancos destruyó información y préstamos, no solo dinero \parencite{Bernanke1983Nonmonetary} \vertema{3.A.37}.\footnote{\autor{Bernanke} recibió el Premio Nobel en 2022, con \autor{Douglas Diamond} y \autor{Philip Dybvig}, «por la investigación sobre bancos y crisis financieras» \parencite{Nobel2022}. Propuesto en 2005, presidió la Reserva Federal de 2006 a 2014.} En el 90 cumpleaños de \autor{Friedman}, siendo gobernador de la Reserva Federal, les dijo a él y a \autor{Schwartz}: «Sobre la Gran Depresión: tenéis razón, lo hicimos nosotros. Lo sentimos mucho. Pero, gracias a vosotros, no lo volveremos a hacer» \parencite[traducción propia]{Bernanke2002Friedman}.\footnote{Lo que vino después, en \href{https://macrothoughts.weebly.com/blog/youre-right-we-did-it-twice}{«“You're right, we did it…”, twice!»} (blog Macro Thoughts).} En 2008, ya como presidente, inundó de liquidez el sistema: esta vez la cantidad de dinero no cayó.
\item \emph{El ciclo monetario}: \textcite{Friedman1964Monetary} describe la producción como una cuerda pegada a un techo, el potencial, de la que de vez en cuando se tira hacia abajo (el \emph{plucking model}), y lo contrasta en \textcite{Friedman1993Plucking}: las recesiones son caídas por debajo del techo, provocadas a menudo por perturbaciones monetarias, y la amplitud de cada caída predice la de la recuperación que la sigue, pero no al revés \vertema{3.A.42}.\footnote{Para ampliar: la explicación del modelo de \autor{Roger Garrison} (\href{https://webhome.auburn.edu/~garriro/fm1pluck.htm}{«Plucking Model»}, Auburn University), \href{https://www.businessinsider.com/milton-friedmans-plucking-model-and-europe-2015-2}{«Milton Friedman's “Plucking Model” and Europe»} (\emph{Business Insider}, 2015) y \autor{Russ Roberts} en \href{https://youtu.be/1eCYq2vD5GY?t=490}{«The Numbers Game: The Economic Recovery»} (Hoover Institution, YouTube).}
\item \emph{Valoración}: \emph{A Monetary History} rehabilitó las explicaciones monetarias del ciclo, que la revolución keynesiana había arrinconado; treinta años después, \textcite{Lucas1994Review} le dedicó una reseña elogiosa.
\end{itemize}

% --------------------------------------------------------------------------
\subsection{Valoración y aspectos criticados (la contracrítica)}\label{sec:valoracion}

\begin{itemize}
\item \emph{Virtudes}: lo que aportó el monetarismo y por qué lo seguimos estudiando:
  \begin{itemize}
  \item la \emph{tasa natural de paro} y la curva de \autor{Phillips} vertical a largo plazo \vertema{3.A.28};
  \item una política monetaria centrada en el \emph{control de la inflación}, responsabilidad del banco central;
  \item las \emph{expectativas} en el centro de los modelos;
  \item las \emph{reglas} frente a la discrecionalidad, con los retardos de la política, una idea que llevó más lejos \autor{Lucas};
  \item los \emph{tipos de cambio flexibles}, que se generalizaron desde 1973.
  \end{itemize}
\item \emph{La contracrítica en la práctica}: las contracciones monetarias de finales de los 70 y principios de los 80 costaron mucho en renta y empleo. En Estados Unidos, el paro del 10,8~\% de 1982 fue el más alto de la posguerra hasta entonces \parencite{BLS2026UNRATE}, aunque también pesó el segundo choque del petróleo. \textcite{Kaldor1970New,Kaldor1982Scourge} atribuyó al monetarismo la recesión británica de 1980-1981 y sostuvo que la oferta monetaria es endógena: la crea el crédito que piden las empresas, y el banco central no la controla.
\item \emph{La contracrítica en la teoría}:
  \begin{itemize}
  \item \emph{causación inversa}: que el dinero se adelante a la renta no prueba que la cause. \textcite{Tobin1970Money} mostró que un modelo keynesiano en el que el dinero responde a la renta esperada genera los mismos adelantos que los monetaristas tomaban por causalidad;
  \item \emph{la cantidad de dinero como objetivo intermedio}: es la \emph{ley de \autor{Goodhart}}: toda regularidad estadística tiende a romperse cuando se la usa para el control \parencite{Goodhart1975Problems}. La relación entre dinero y renta se deshizo justo cuando los bancos centrales se apoyaron en ella; si las perturbaciones vienen de la demanda de dinero, estabiliza mejor fijar el tipo de interés \parencite{Poole1970Optimal}. Los bancos centrales pasaron a reglas de tipo de interés, como la de \autor{Taylor} \vertema{3.A.36};
  \item \emph{la renta permanente} no se observa, y es difícil de contrastar;
  \item \emph{las expectativas adaptativas} permiten errores sistemáticos: la nueva macroeconomía clásica las sustituye por las racionales de \textcite{Muth1961Rational} y llega a una conclusión aún más dura contra la estabilización \parencite{Lucas1972Expectations} \vertema{3.A.6}.
  \end{itemize}
\item \emph{El legado}: la nueva síntesis neoclásica de los 90, la nueva economía keynesiana con expectativas racionales, recoge buena parte de las ideas monetaristas (tasa natural, expectativas, reglas, inflación como objetivo) con las rigideces de la síntesis \parencite{GoodfriendKing1997New} \vertema{3.A.7}. Como resume \textcite{Mankiw2006Macroeconomist}, esa nueva síntesis es la que usan hoy los bancos centrales.
\item \emph{La paradoja final}: \autor{Friedman} cambió los términos del debate y puso a los keynesianos a la defensiva, pero lo que vino después no fue el monetarismo, sino la nueva macroeconomía clásica, walrasiana y con expectativas racionales. Muchos monetaristas descubrieron entonces que, marshallianos como sus viejos adversarios keynesianos, tenían más en común con ellos que con los nuevos clásicos \parencite{DeVroey2016History}.
\end{itemize}

\begin{table}[H]
\centering
\caption{La síntesis neoclásica y el monetarismo, pieza a pieza}\label{tab:sintesis-monetarismo}
\small
\begin{tabular}{>{\raggedright\arraybackslash}m{3.1cm}>{\raggedright\arraybackslash}m{5.5cm}>{\raggedright\arraybackslash}m{5.5cm}}
\toprule
\rowcolor{tceeazulpalido} & \textcolor{tceemorado}{\textbf{Síntesis neoclásica}} & \textcolor{tceeazul}{\textbf{Monetarismo}}\\
\midrule
\textbf{Consumo} & \textcolor{tceemorado}{$C_t=\overline{C}+c\,Y^d_t$: renta corriente} & \textcolor{tceeazul}{$C^P_t\propto Y^P_t$: renta permanente}\\
\textbf{Inversión} & \textcolor{tceemorado}{poco sensible al tipo: IS inclinada} & \textcolor{tceeazul}{gasto amplio y sensible: IS plana}\\
\textbf{Demanda de dinero} & \textcolor{tceemorado}{inestable, $h$ grande: LM plana} & \textcolor{tceeazul}{estable, $h$ pequeño: LM inclinada}\\
\textbf{Transmisión} & \textcolor{tceemorado}{tipo de interés de los bonos} & \textcolor{tceeazul}{reajuste de toda la cartera}\\
\textbf{Expectativas} & \textcolor{tceemorado}{estáticas o exógenas} & \textcolor{tceeazul}{adaptativas}\\
\textbf{Curva de \autor{Phillips}} & \textcolor{tceemorado}{menú permanente} & \textcolor{tceeazul}{vertical a largo plazo en $U^n$}\\
\textbf{Economía privada} & \textcolor{tceemorado}{inestable} & \textcolor{tceeazul}{estable}\\
\textbf{Política monetaria} & \textcolor{tceemorado}{débil y secundaria; discrecional, con el tipo de interés} & \textcolor{tceeazul}{potente, con retardos; regla del $k$~\% sobre la cantidad de dinero}\\
\textbf{Política fiscal} & \textcolor{tceemorado}{preferida, discrecional (sintonía fina); poca expulsión} & \textcolor{tceeazul}{mucha expulsión; reglas y estabilizadores automáticos}\\
\textbf{Políticas de oferta} & \textcolor{tceemorado}{políticas de rentas} & \textcolor{tceeazul}{reformas que bajan $U^n$; sin controles de precios}\\
\textbf{Tipo de cambio} & \textcolor{tceemorado}{fijo y ajustable (Bretton Woods)} & \textcolor{tceeazul}{flexible}\\
\textbf{Balanza de pagos} & \textcolor{tceemorado}{elasticidades y absorción} & \textcolor{tceeazul}{fenómeno monetario (enfoque monetario)}\\
\textbf{Gran Depresión} & \textcolor{tceemorado}{demanda insuficiente} & \textcolor{tceeazul}{contracción monetaria}\\
\bottomrule
\end{tabular}
\fuente{elaboración propia.}
\end{table}

\begin{itemize}
\item \emph{Recapitulación}: el monetarismo convirtió su crítica en una receta (reglas y no sintonía fina, reformas de oferta y tipos flexibles) que dominó la política a principios de los 80 y se rompió con la velocidad del dinero. Lo que sobrevivió, la tasa natural, las expectativas y las reglas, pasó a la macroeconomía posterior.
\end{itemize}


% ==========================================================================
\section*{Conclusión}
\minutos{1}

\begin{itemize}
\item \textbf{\emph{Recapitulación}}:
  \begin{itemize}
  \item En conclusión, la síntesis neoclásica convirtió a \autor{Keynes} en un modelo, el IS-LM, cerrado con la curva de \autor{Phillips}: \autor{Keynes} a corto plazo, los neoclásicos a largo y una política de demanda que elegía en un menú de inflación y paro.

  \item El monetarismo mantuvo el modelo, pero cambió sus piezas: un consumo y una demanda de dinero estables, un dinero potente y una curva de \autor{Phillips} vertical a largo plazo.

  \item De ahí salió su receta: una regla monetaria del $k$~\%, una política fiscal sin sintonía fina, reformas que bajen la tasa natural y tipos de cambio flexibles.

  \item \autor{Harry Johnson} lo llamó «contrarrevolución» monetarista \parencite{Johnson1971Keynesian}, una expresión que el propio \textcite{Friedman1970Counter} había usado un año antes.
  \end{itemize}
\item \textbf{\emph{Relevancia}}:
  \begin{itemize}
  \item La estanflación de los 70 dio la razón a la crítica, y la desinflación de 1979-1982 y la inestabilidad de la velocidad, a la contracrítica.

  \item Lo que sobrevivió es el suelo de los bancos centrales de hoy.
  \end{itemize}
\item \textbf{\emph{Extensiones}}:
  \begin{itemize}
  \item El desequilibrio \vertema{3.A.4}; las expectativas racionales y la crítica de \autor{Lucas} \vertema{3.A.6}; la nueva economía keynesiana, que microfundamenta las rigideces de la síntesis \vertema{3.A.7}; la demanda de dinero \vertema{3.A.35}; el diseño de la política monetaria, de la regla del $k$~\% a la de \autor{Taylor} \vertema{3.A.36}; la financiación del déficit \vertema{3.A.39}; los regímenes cambiarios \vertema{3.B.15}, y la tasa natural y la NAIRU \vertema{3.A.28}.
  \end{itemize}
\item \textbf{\emph{Opinión}}:
  \begin{itemize}
  \item Las 2 escuelas tenían parte de razón.

  \item La síntesis acertó en que la demanda importa a corto plazo y en que la política puede estabilizar; el monetarismo, en que no hay atajos permanentes contra el paro y en que el dinero mal gestionado causa las peores crisis.

  \item Se equivocó en la herramienta: la cantidad de dinero resultó un mal objetivo.
  \end{itemize}
\item \textbf{\emph{Idea final}}:
  \begin{itemize}
  \item «Cada maestrillo tiene su librillo», pero el libro era el mismo.

  \item Keynesianos y monetaristas compartían el IS-LM y la curva de \autor{Phillips}; lo que los separaba eran las pendientes, la estabilidad del dinero y las expectativas.

  \item Por eso, como dijo \autor{Friedman}, en un sentido todos somos hoy keynesianos y, en otro, nadie lo es ya: los bancos centrales leen el mismo libro con un librillo que ha tomado páginas de los 2.
  \end{itemize}
\end{itemize}

\clearpage
\appendix

% ==========================================================================
\section{Anexo: el modelo IS-LM paso a paso}\label{anexo:islm}

\begin{itemize}
\item \emph{Las 2 curvas}. De la IS \eqref{eq:is} y la LM \eqref{eq:lm},
  \begin{align}
  \big[1-c\,(1-t)\big]\,Y&=\overline{A}-b\,i, \label{eq:anexo-is}\\
  i&=\frac{1}{h}\Big(k\,Y-\frac{M}{P}\Big). \label{eq:anexo-lm}
  \end{align}
\item \emph{El equilibrio}. Sustituyendo \eqref{eq:anexo-lm} en \eqref{eq:anexo-is} y multiplicando por $h$,
  \begin{equation*}
  \big\{h\,\big[1-c\,(1-t)\big]+b\,k\big\}\,Y=h\,\overline{A}+b\,\frac{M}{P}.
  \end{equation*}
  Con $D\equiv h\,[1-c\,(1-t)]+b\,k>0$ (local de este anexo),
  \begin{equation*}
  Y^*=\frac{h\,\overline{A}+b\,M/P}{D},\qquad i^*=\frac{k\,Y^*-M/P}{h}=\frac{k\,\overline{A}-[1-c\,(1-t)]\,M/P}{D},
  \end{equation*}
  que es \eqref{eq:equilibrio-islm}.
\item \emph{Multiplicadores}:
  \begin{equation*}
  \frac{\partial Y^*}{\partial\overline{G}}=\frac{h}{D}=\frac{1}{1-c\,(1-t)+b\,k/h},\qquad \frac{\partial i^*}{\partial\overline{G}}=\frac{k}{D}>0,
  \end{equation*}
  \begin{equation*}
  \frac{\partial Y^*}{\partial(M/P)}=\frac{b}{D},\qquad \frac{\partial i^*}{\partial(M/P)}=-\frac{1-c\,(1-t)}{D}<0.
  \end{equation*}
\item \emph{El efecto expulsión}. La inversión cae $\Delta I=-b\,\Delta i^*=-(b\,k/D)\,\Delta\overline{G}$. La renta crece lo que dice el multiplicador simple aplicado al gasto neto de la expulsión:
  \begin{equation*}
  \frac{\Delta\overline{G}+\Delta I}{1-c\,(1-t)}=\frac{\Delta\overline{G}\,(D-b\,k)}{D\,[1-c\,(1-t)]}=\frac{h}{D}\,\Delta\overline{G},
  \end{equation*}
  porque $D-b\,k=h\,[1-c\,(1-t)]$.
\item \emph{Los casos extremos} (tabla~\ref{tab:casos-islm}), tomando límites:
  \begin{itemize}
  \item $h\to\infty$: $h/D\to 1/[1-c\,(1-t)]$ y $b/D\to 0$ (trampa de la liquidez);
  \item $h=0$: $D=b\,k$, $h/D=0$ y $b/D=1/k$; es decir, $Y=(1/k)\,M/P$, la teoría cuantitativa con $V=1/k$ (caso clásico);
  \item $b=0$: $D=h\,[1-c\,(1-t)]$, $h/D=1/[1-c\,(1-t)]$ y $b/D=0$ (inversión insensible);
  \item $b\to\infty$: $h/D\to 0$ y $b/D\to 1/k$ (inversión muy sensible).
  \end{itemize}
\item \emph{La demanda agregada}. $Y^*$ en función de $P$ es la DA \eqref{eq:da}, con pendiente
  \begin{equation*}
  \frac{\partial Y^*}{\partial P}=-\frac{b}{D}\,\frac{M}{P^2}<0,
  \end{equation*}
  que se anula si $b=0$ o si $h\to\infty$: la DA es vertical. Con el efecto saldos reales, el gasto autónomo depende de los saldos reales, $\overline{A}(M/P)$ con $\overline{A}'>0$, y
  \begin{equation*}
  \frac{\partial Y^*}{\partial P}=-\frac{h\,\overline{A}'+b}{D}\,\frac{M}{P^2}<0
  \end{equation*}
  aun con $b=0$, y con $h\to\infty$ tiende a $-\overline{A}'\,M/\{[1-c\,(1-t)]\,P^2\}<0$: la deflación siempre acaba elevando la demanda (efecto \autor{Pigou}-\autor{Patinkin}).
\end{itemize}
\findemo

% ==========================================================================
\section{Anexo: la metodología de \autor{Friedman}}\label{anexo:metodo}

\begin{itemize}
\item \emph{La predicción como criterio}: para \textcite{Friedman1953Methodology}, la economía positiva es una ciencia que debe abordar problemas concretos y dar predicciones que se puedan contrastar con los datos; una teoría se acepta mientras no la refuten. Es afín al falsacionismo de \autor{Popper} (una teoría no se puede verificar, pero sí refutar con una observación incompatible) \vertema{3.A.1}, aunque con un instrumentalismo que \autor{Popper} no compartía: las teorías son instrumentos para predecir, no descripciones verdaderas del mundo \parencite{Boland1979Critique}.
\item \emph{Economía positiva frente a normativa}: \autor{Friedman} defiende que la economía positiva puede ser independiente de cualquier postura ética y que la normativa debe apoyarse en ella. Frente a él, \autor{Gunnar Myrdal} sostenía que los conceptos económicos llevan siempre valoraciones implícitas, que hay que hacer explícitas \parencite{Myrdal1953Political}.\footnote{\autor{Myrdal} recibió el Premio Nobel en 1974, con \autor{Friedrich Hayek}, «por sus trabajos pioneros en la teoría del dinero y de las fluctuaciones económicas y por su penetrante análisis de la interdependencia de los fenómenos económicos, sociales e institucionales» \parencite{Nobel1974}.} Las recomendaciones de \autor{Friedman} se presentan así: controlar el crecimiento del dinero se deriva de un resultado positivo, que el dinero es lo que más explica la inflación.
\item \emph{Irrelevancia del realismo de los supuestos}: lo que cuenta es lo bien que predice una teoría, no lo realistas que sean sus supuestos. Sus ejemplos: el buen jugador de billar golpea \emph{como si} resolviera las ecuaciones de la mecánica, y las empresas se comportan \emph{como si} maximizaran el beneficio, aunque no calculen ingresos ni costes marginales: las que no lo hacen desaparecen. \autor{Samuelson} lo criticó como el «giro F» (\emph{F-twist}): unos supuestos falsos no se vuelven inocuos porque el modelo prediga bien \parencite{Samuelson1963Problems}.
\item \emph{Sencillez}: una buena teoría explica mucho con poco; \autor{Friedman} prefiere los modelos pequeños y el razonamiento económico verbal al detalle técnico.
\item \emph{Enfoque marshalliano}: la economía como «máquina para descubrir verdades concretas», que estudia cada problema por separado, con equilibrio parcial y datos, frente al equilibrio general walrasiano, que veía como un ejercicio formal \parencite{Friedman1949Marshallian}. Paradójicamente, \autor{Brunner} y \autor{Meltzer}, sus aliados, trabajaban en la tradición walrasiana \parencite{BordoSchwartz2003ISLM}.
\end{itemize}

% ==========================================================================
\section{Anexo: aspectos para el test}\label{anexo:test}

\begin{itemize}
\item \emph{Fechas y obras}:
  \begin{itemize}
  \item \autor{Keynes}, \emph{Teoría general} (1936); \autor{Hicks}, «Mr.~Keynes and the ``Classics''; A Suggested Interpretation» (1937, curvas IS y LL); \autor{Modigliani} (1944, salarios rígidos); \autor{Hansen}, \emph{Monetary Theory and Fiscal Policy} (1949) y \emph{A Guide to Keynes} (1953); \autor{Samuelson}, \emph{Economics} (1.ª ed., 1948; «síntesis neoclásica» en la 3.ª, 1955); \autor{Klein} y \autor{Goldberger} (1955); \autor{Phillips} (1958); \autor{Samuelson} y \autor{Solow} (1960);
  \item \autor{Friedman}: «The Methodology of Positive Economics» y «The Case for Flexible Exchange Rates» (1953), «The Quantity Theory of Money: A Restatement» (1956), \emph{A Theory of the Consumption Function} (1957), \emph{A Program for Monetary Stability} (1960), \emph{Capitalism and Freedom} (1962), «The Role of Monetary Policy» (1968, discurso de diciembre de 1967) y \emph{The Optimum Quantity of Money} (1969); \autor{Friedman} y \autor{Schwartz}, \emph{A Monetary History of the United States, 1867-1960} (1963);
  \item \autor{Cagan} (1956) y \autor{Nerlove} (1958), expectativas adaptativas; \autor{Phelps} (1967, 1968); \autor{Brunner} (1968), el nombre «monetarismo»; \autor{Johnson} (1971), «contrarrevolución»; \autor{Tobin} (1970), causación inversa; \autor{Poole} (1970); \autor{Goodhart} (1975).
  \end{itemize}
\item \emph{Nobel}: \autor{Friedman} en 1976; \autor{Myrdal} y \autor{Hayek} en 1974 («interdependencia» de los fenómenos económicos, sociales e institucionales); \autor{Bernanke} en 2022 \parencite{Nobel1976Friedman,Nobel1974,Nobel2022}.
\item \emph{Trampas}:
  \begin{itemize}
  \item \autor{Hicks} llamó LL a su curva monetaria; el nombre IS-LM se generalizó con \autor{Hansen};
  \item la trampa de la liquidez es la LM \emph{horizontal} ($h\to\infty$); el caso clásico, la LM \emph{vertical} ($h=0$);
  \item los monetaristas: IS \emph{plana} y LM \emph{inclinada}; los keynesianos, al revés;
  \item la regla del $k$~\% (1960) no es la «regla de \autor{Friedman}» de la cantidad óptima de dinero (1969);
  \item \autor{Keynes} no supuso constante la propensión marginal a consumir: fue la síntesis;
  \item el efecto \autor{Keynes} actúa por el tipo de interés (desplaza la LM); el de \autor{Pigou}-\autor{Patinkin}, por la riqueza (desplaza la IS);
  \item la curva de \autor{Phillips} de 1958 relaciona el paro con los \emph{salarios} nominales; la de \autor{Samuelson} y \autor{Solow}, con los \emph{precios};
  \item en la Gran Contracción la base monetaria \emph{subió}; lo que cayó fue el multiplicador.
  \item la política fiscal de \autor{Friedman} (1948): estabilizadores automáticos y presupuesto equilibrado en pleno empleo; la regla del $k$~\% llega en 1960;
  \item el enfoque monetario de la balanza de pagos (\autor{Harry Johnson}): con tipo fijo, un déficit exterior es un exceso de oferta de dinero que se corrige con la pérdida de reservas, $\Delta R=\Delta M^D-\Delta D$;
  \item los monetaristas rechazan las políticas de rentas (controles de precios y salarios) y prefieren la desinflación gradual; la de golpe es de la nueva macroeconomía clásica.
  \end{itemize}
\item Preguntas de test de exámenes anteriores, en el simulador: \url{https://www.victorgutierrezmarcos.es/oposicion/temario/primer-ejercicio/test/simulador.html}.
\end{itemize}

\bibliografia

\end{document}
